%% Wiki: [[2026.01.22 - Mathematics: Scheduling and Compiler Techniques for Metadata-Hiding Overlay Lookups]]
\documentclass[11pt]{article}

\usepackage[margin=1in]{geometry}
\usepackage[T1]{fontenc}
\usepackage[utf8]{inputenc}
\usepackage{lmodern}
\usepackage{amsmath,amssymb,mathtools}
\usepackage{microtype}
\usepackage{booktabs,longtable,array}
\usepackage{enumitem}
\usepackage[hidelinks]{hyperref}
\usepackage{newunicodechar}

% Pandoc compatibility
\providecommand{\tightlist}{%
  \setlength{\itemsep}{0pt}\setlength{\parskip}{0pt}}

% Unicode shims (pdflatex-friendly).
\newunicodechar{ε}{\ensuremath{\varepsilon}}
\newunicodechar{τ}{\ensuremath{\tau}}
\newunicodechar{β}{\ensuremath{\beta}}
\newunicodechar{δ}{\ensuremath{\delta}}
\newunicodechar{μ}{\ensuremath{\mu}}
\newunicodechar{ν}{\ensuremath{\nu}}
\newunicodechar{π}{\ensuremath{\pi}}
\newunicodechar{Δ}{\ensuremath{\Delta}}
\newunicodechar{Θ}{\ensuremath{\Theta}}
\newunicodechar{Φ}{\ensuremath{\Phi}}
\newunicodechar{≤}{\ensuremath{\le}}
\newunicodechar{≥}{\ensuremath{\ge}}
\newunicodechar{≈}{\ensuremath{\approx}}
\newunicodechar{∞}{\ensuremath{\infty}}
\newunicodechar{∪}{\ensuremath{\cup}}
\newunicodechar{→}{\ensuremath{\to}}
\newunicodechar{⇒}{\ensuremath{\Rightarrow}}
\newunicodechar{−}{-}
\newunicodechar{–}{--}
\newunicodechar{’}{'}
\newunicodechar{“}{``}
\newunicodechar{”}{''}
\newunicodechar{·}{\ensuremath{\cdot}}
\newunicodechar{≳}{\ensuremath{\gtrsim}}
\newunicodechar{□}{\ensuremath{\square}}

\setlist[itemize]{leftmargin=1.2em}
\setlist[enumerate]{leftmargin=1.5em}

\title{Scheduling and Compiler Techniques for Metadata-Hiding Overlay Lookups}
\author{}
\date{Draft v0.67 (2026-01-21). Anonymous submission; author list
omitted.}

\begin{document}
\maketitle

\noindent\textbf{Canonical wiki page:} \texttt{[[2026.01.22 - Mathematics: Scheduling and Compiler Techniques for Metadata-Hiding Overlay Lookups]]}\\


\hypertarget{abstract}{%
\section{Abstract}\label{abstract}}

We study two metadata channels that persist even when payload privacy is
perfect: \textbf{(i)} observable \emph{termination / stop time}, and
\textbf{(ii)} observable \emph{destination / contact traces} over
repeated operations. We develop a proof-oriented \textbf{compiler
calculus}: small, auditable transformations that export quantitative
contracts (TV, pointwise, KL/MI) and can be composed into end-to-end
leakage guarantees.

For stop times, we analyze a simple ε-tunable \textbf{mixture compiler}
and show it is optimal \emph{within the mixture class}. We then show
that broader \textbf{kernel-feasible} padding rules can be strictly
better: we formulate globally optimal kernel-feasible padding as a
linear program (\textbf{Kernel-TV-Pad}), give explicit separating
instances and parametric families where mixtures are arbitrarily
suboptimal, and provide machine-checkable primal--dual certificates.

For contact traces, we combine \textbf{public scheduling} (PSC), which
makes the set of coarse destinations a deterministic function of public
time, with \textbf{mixture equalization} (MC-EQ), which equalizes within
each scheduled set via one-real-among-w mixing. Our main composition
theorems convert these module contracts into (i)
transcript-by-transcript \textbf{posterior} guarantees near \(1/K\) and
(ii) \textbf{MI/Fano} identification bounds, together with a
reproducible parameter-selection worksheet and figures.

\begin{center}\rule{0.5\linewidth}{0.5pt}\end{center}

\hypertarget{motivation-leakage-channels-and-notation}{%
\section{1. Motivation, leakage channels, and
notation}\label{motivation-leakage-channels-and-notation}}

Overlay lookup protocols (including many DHT-like designs) can hide
\emph{what} is being queried while still leaking \emph{how the query
unfolds}. This paper treats two particularly stubborn metadata channels:
- \textbf{Stop time}: when the observable protocol terminates. -
\textbf{Contact trace / marks}: which coarse destinations (e.g.,
shards/committees) are touched at each time slot. Our goal is a
\textbf{compiler calculus}: small, auditable transformations that (i)
fit implementability constraints (causality, feasibility), and (ii)
admit quantitative leakage bounds that compose. We emphasize \textbf{two
complementary leakage regimes}: - \textbf{Posterior-first} (near-\(1/K\)
posteriors): protect against rare transcripts that spike the posterior.
- \textbf{MI/Fano} (average information): useful when only
high-confidence identification needs to be ruled out. A transcript
consists of: - a stopping time \(T\) (termination-only leakage), or - a
contact-mark sequence \(O_{1:T}\) (e.g., ``which shard/committee was
contacted in slot \(t\)''), together with a public schedule. We write
\(\mathrm{TV}(P,Q)\) for total-variation distance and
\(D_{\mathrm{KL}}(P\|Q)\) for KL divergence (nats). We use \(I(C;O)\)
for MI between a hidden client/profile variable \(C\) and observations
\(O\). ---

\hypertarget{a-compiler-calculus-interface-modules-and-contracts}{%
\subsection{1.1 A compiler-calculus interface (modules and
contracts)}\label{a-compiler-calculus-interface-modules-and-contracts}}

To keep the end-to-end arguments modular, we treat each transformation
as exporting a small \textbf{contract} that can be consumed by a
composition theorem. Informally: - the \textbf{trace compiler} exports
per-slot deviation bounds relative to a secret-independent reference
kernel, - the \textbf{stop-time compiler} exports a TV /
posterior-friendly bound for the padded termination distribution, - the
\textbf{composition theorem} consumes these bounds to control either
worst-case \textbf{posteriors} or average \textbf{MI/Fano} risk. We
summarize the main modules and their ``typed'' outputs in Table 1.
\textbf{Table 1 (Module contracts used by the composition theorems).}

\begin{longtable}[]{@{}
  >{\raggedright\arraybackslash}p{(\columnwidth - 6\tabcolsep) * \real{0.2500}}
  >{\raggedright\arraybackslash}p{(\columnwidth - 6\tabcolsep) * \real{0.2500}}
  >{\raggedright\arraybackslash}p{(\columnwidth - 6\tabcolsep) * \real{0.2500}}
  >{\raggedright\arraybackslash}p{(\columnwidth - 6\tabcolsep) * \real{0.2500}}@{}}
\toprule
\begin{minipage}[b]{\linewidth}\raggedright
Module
\end{minipage} & \begin{minipage}[b]{\linewidth}\raggedright
What it does
\end{minipage} & \begin{minipage}[b]{\linewidth}\raggedright
Exported contract (parameters)
\end{minipage} & \begin{minipage}[b]{\linewidth}\raggedright
Consumed by
\end{minipage} \\
\midrule
\endhead
\textbf{PSC (public scheduling)} & Makes the action-set schedule \(J_t\)
a deterministic function of public time; collapses shard/committee marks
to within-set choice. & Public sets \(J_t\); (optionally) a
support-containment guarantee for KL/MI accounting (Remark 3.1). &
Sections 3--6; Theorems 10.1--10.2. \\
\textbf{MC-EQ (mixture equalization)} & Within a scheduled set \(J_t\)
of size \(m\), mixes one secret-dependent action with \(w-1\) covers
from a secret-independent kernel \(K_t(\cdot\mid h)\). & Per-slot
deviations relative to \(K_t\): \(\mathrm{TV}\le 1/w\) and pointwise
\(|P(y\mid c,h)-K_t(y\mid h)|\le 1/w\) (Lemma 4.2). & Theorem 7.2;
Theorems 10.1--10.2. \\
\textbf{StopTimePad (Kernel-TV-Pad)} & Kernel-feasible padding
\(\Phi(t\mid s)\) supported on \(t\ge s\) that minimizes expected padded
stop time under a pairwise TV budget. & Pairwise stop-time privacy:
\(\mathrm{TV}(P_c,P_{c'})\le \tau\) for all \(c,c'\). & Section 2;
Theorem 10.2 (MI/Fano path). \\
\textbf{StopTimePad (PosteriorStopTime-Pad)} & Stop-time padding tuned
for posterior-first composition: enforces an \(\ell_\infty\)-star
constraint around a reference \(R\) with a probability floor. &
Existence of \(R\) with floor \(\kappa_T\) and additive deviations
\(\delta_T\): \(|P_c(t)-R(t)|\le \delta_T\), \(R(t)\ge \kappa_T\). &
Proposition 2.15; Theorem 10.1 (posterior-first path). \\
\textbf{Composition} & Converts per-slot contracts (trace + stop time)
into transcript-level leakage bounds. & (i) Posterior guarantee (Theorem
10.1), (ii) MI/KL bound for Fano (Theorem 10.2). & End-to-end results +
design worksheet (Section 10.4). \\
\bottomrule
\end{longtable}

This interface perspective is what lets Option B (``compiler calculus'')
scale: adding a new compiler is useful whenever it exports a contract
that composes.

\hypertarget{threat-model-and-scope}{%
\subsection{1.2 Threat model and scope}\label{threat-model-and-scope}}

We assume an adversary who observes the \textbf{public schedule} (PSC),
the per-slot \textbf{contact marks} (e.g., shard/committee identifiers
contacted in each slot), and the final \textbf{stop time}. Payload
contents are protected by standard cryptography (e.g., PIR inside a
shard, encryption, authentication), and real versus dummy messages are
assumed computationally indistinguishable at the payload layer. Our
focus is on \emph{metadata} that remains: the observable sequence of
marks and the termination time. \textbf{What we cover.} - Coarse
destination leakage (marks) at the granularity of scheduled action sets.
- Stop-time leakage under implementable (causal, monotone) padding
rules. - Two leakage regimes: worst-case \textbf{posterior spikes}
(near-\(1/K\) posteriors) and average \textbf{MI/Fano} leakage.
\textbf{What we do not cover.} We do not model packet-level timing
within a slot, size/shape leakage, routing-layer side channels,
churn-induced linkability, or active attacks that break the
payload-layer indistinguishability assumption. These can be layered on
top, but are orthogonal to the compiler calculus developed here.

\hypertarget{main-results-at-a-glance}{%
\subsection{1.3 Main results at a
glance}\label{main-results-at-a-glance}}

The paper is organized around a modular pipeline; the most important
results can be read as \emph{typing rules} for modules and as
\emph{composition theorems} that consume their contracts.
\textbf{Stop-time compilers (Section 2).} - Mixture compilers achieve
StopTime-IND\(_\varepsilon\) with a linear privacy--overhead dial and
are tight \emph{within the mixture class}. - Kernel-feasible padding is
strictly more powerful: we give explicit separations, an LP formulation
(\textbf{Kernel-TV-Pad}) with a dual/KKT form (certificate-friendly),
and a closed-form parametric family where mixtures are
\textbf{arbitrarily suboptimal} (gap \(\Omega(n)\)); see Proposition 2.2
and Theorem 2.8. - For \(K\)-ary stop-time hiding, we formalize
\textbf{pairwise} vs \textbf{TV-star} constraints, prove a tight TV
geometry lemma (Lemma 2.13), and provide machine-checkable primal--dual
certificates for strict gaps (Section 2.9). \textbf{Trace compilers and
composition (Sections 3--10).} - PSC collapses shard-level marks to
within-set choices. - MC-EQ provides a tunable equalization knob \(w\)
and exports both TV and pointwise deviation contracts (Lemma 4.2). - The
flagship \textbf{compiler calculus theorems} (Theorems 10.1 and 10.2)
compose PSC + MC-EQ + stop-time padding into (i)
transcript-by-transcript posterior guarantees and (ii) MI/Fano
guarantees, with a worked parameter-selection worksheet (Section 10.4).
\textbf{Figure 1 (pipeline).} A one-page diagram in Appendix A
summarizes the module contracts and how they feed into Theorems
10.1--10.2.

\hypertarget{notation-quick-reference}{%
\subsection{1.4 Notation quick
reference}\label{notation-quick-reference}}

We use the following symbols repeatedly; the goal is to keep the later
proofs ``type-checkable'' at a glance. - \(C\): hidden
client/profile/secret, typically uniform on \([K]\). - \(K\): number of
possible secrets/profiles. - \(T\): (padded) observable stop time;
\(S\): (un-padded) success time. - \(n\): discrete time horizon for
stop-time padding examples. - \(O_t\): observable mark in slot \(t\)
(e.g., which action within the scheduled set was contacted). - \(J_t\):
scheduled action set in slot \(t\) (public under PSC), with \(|J_t|=m\).
- \(m\): scheduled-set size (number of actions available in a slot). -
\(w\): equalization fanout / cover multiplicity (``one real among
\(w\)''), with \(\rho:= 1/w\). - \(\tau\): TV privacy budget for
stop-time distributions (pairwise \(\mathrm{TV}(P_c,P_{c'})\le \tau\)).
- \(\delta\): generic per-step deviation parameter (TV or pointwise),
with floors \(\kappa\) when converting to KL. - \(\mathrm{TV}(P,Q)\):
total variation; \(D_{\mathrm{KL}}(P\|Q)\): KL divergence (nats);
\(I(C;O)\): mutual information.

\hypertarget{reproducibility}{%
\subsection{1.5 Reproducibility}\label{reproducibility}}

This draft is accompanied by small deterministic scripts that reproduce
every numerical claim and primal--dual certificate. To run the full
suite, execute:

\begin{verbatim}
bash reproduce_all_v67.sh
\end{verbatim}

Run the command from the directory that contains the scripts (e.g., the
artifact root). The runner is path-robust and will regenerate Figures
2--3.

Typical runtime on a laptop is under a minute; most certificates
complete in seconds.

Appendix C lists the mapping from claims to scripts.

\hypertarget{how-to-read-this-paper}{%
\subsection{1.6 How to read this paper}\label{how-to-read-this-paper}}

\begin{itemize}
\tightlist
\item
  For \textbf{stop-time padding} (timing-only leakage), start with
  Section 2. Sections 2.4--2.10 contain the LP formulations, separations
  from mixture compilers, and primal--dual certificates.
\item
  For \textbf{contact traces} (who is contacted when), read Sections
  3--6: PSC collapses shard-mark leakage, MC-EQ provides the per-slot
  equalization knob, and we relate per-step mixing to the hypotheses of
  the sequential theorems.
\item
  For an \textbf{end-to-end privacy guarantee}, jump to Theorem 10.1
  (posterior-first, near-\(1/K\) posteriors) and Theorem 10.2 (MI/Fano,
  high-confidence identification risk). Section 10.4 gives a worked
  design worksheet.
\end{itemize}

\begin{center}\rule{0.5\linewidth}{0.5pt}\end{center}

\hypertarget{stop-time-leakage-stoptime-ind-and-mixture-compilation}{%
\section{2. Stop-time leakage: StopTime-IND and mixture
compilation}\label{stop-time-leakage-stoptime-ind-and-mixture-compilation}}

\hypertarget{stoptime-ind_varepsilon}{%
\subsection{\texorpdfstring{2.1
StopTime-IND\(_\varepsilon\)}{2.1 StopTime-IND\_\textbackslash varepsilon}}\label{stoptime-ind_varepsilon}}

In the stop-time-only leakage model, the adversary observes only the
stop time \(T\). StopTime-IND\(_\varepsilon\) is equivalent to
\(\|P_0-P_1\|_{\mathrm{TV}}\le 2\varepsilon\).

\hypertarget{mixture-compiler-and-tightness-within-the-mixture-class}{%
\subsection{2.2 Mixture compiler and tightness within the mixture
class}\label{mixture-compiler-and-tightness-within-the-mixture-class}}

A mixture compiler \(C_{\lambda,Q}\) (probability \(\lambda\) stop at
true success time \(S\); otherwise stop at secret-independent
\(T_Q\sim Q\) with \(T_Q\ge S\) a.s.) yields
\(\|P_0-P_1\|_{\mathrm{TV}}\le \lambda\), hence choosing
\(\lambda=2\varepsilon\) gives StopTime-IND\(_\varepsilon\).
\textbf{Lemma 2.1 (tightness within the timeout-mixture class).}
Consider the following timeout-mixture padding rule: fix a
secret-independent distribution \(Q\) supported on times \(t\) that are
always feasible (e.g., a deterministic public timeout at the horizon
\(n\)), and output \(T=S\) with probability \(\lambda\) and \(T\sim Q\)
with probability \(1-\lambda\). Then for any two worlds \(b,b'\), \[
\mathrm{TV}(P_b,P_{b'}) = \lambda\,\mathrm{TV}(S_b,S_{b'}),
\] so in particular \(\mathrm{TV}(P_b,P_{b'})\le \lambda\) and this
bound is \textbf{tight} whenever \(\mathrm{TV}(S_b,S_{b'})=1\).
Consequently, to satisfy StopTime-IND\(_\varepsilon\) against worst-case
pairs in this class one must have \(\lambda\le 2\varepsilon\), and
choosing \(\lambda=2\varepsilon\) is optimal.

\emph{Proof.} By construction \(P_b=(1-\lambda)Q+\lambda S_b\), so
\(P_b-P_{b'}=\lambda(S_b-S_{b'})\) and TV scales linearly. \(\square\)

\hypertarget{global-optimal-stop-time-padding-remains-open}{%
\subsection{2.3 Global optimal stop-time padding remains
open}\label{global-optimal-stop-time-padding-remains-open}}

\hypertarget{beyond-the-timeout-mixture-class-the-design-space-includes-general-kernel-feasible-paddings-section-2.4-and-more-broadly-onlinestateful-padding-rules.-characterizing-globally-optimal-padding-under-additional-constraints-unknown-or-drifting-success-time-laws-continuous-time-tail-constraints-or-adversarial-environments-remains-an-open-direction.}{%
\section{Beyond the timeout-mixture class, the design space includes
general kernel-feasible paddings (Section 2.4) and, more broadly,
online/stateful padding rules. Characterizing globally optimal padding
under additional constraints (unknown or drifting success-time laws,
continuous time, tail constraints, or adversarial environments) remains
an open
direction.}\label{beyond-the-timeout-mixture-class-the-design-space-includes-general-kernel-feasible-paddings-section-2.4-and-more-broadly-onlinestateful-padding-rules.-characterizing-globally-optimal-padding-under-additional-constraints-unknown-or-drifting-success-time-laws-continuous-time-tail-constraints-or-adversarial-environments-remains-an-open-direction.}}

\hypertarget{kernel-feasible-padding-can-beat-mixture-compilers-new}{%
\subsection{2.4 Kernel-feasible padding can beat mixture compilers
(new)}\label{kernel-feasible-padding-can-beat-mixture-compilers-new}}

The mixture compiler of Section 2.2 is optimal \emph{within its mixture
class}, but it does not exhaust all implementable stop-time padding
rules. In discrete time, an implementable padding rule can be modeled as
a Markov kernel \(\Phi(t\mid s)\) supported on \(t\ge s\), where \(s\)
is the (local) success time. \textbf{Definition 2.1 (kernel-feasible
padding).} Fix a finite horizon \(\{1,\dots,n\}\). A padding kernel is a
stochastic matrix \(\Phi\) with entries \(\Phi(t\mid s)\ge 0\)
satisfying \[
\sum_{t=s}^n \Phi(t\mid s)=1\quad\text{and}\quad \Phi(t\mid s)=0\ \text{for}\ t<s.
\] Given a success time \(S=s\), the padded stop time \(T\) is sampled
as \(T\sim\Phi(\cdot\mid s)\). For world \(b\in\{0,1\}\) with
success-time law \(S_b\), the induced stop-time law is \[
P_b(t)=\sum_{s=1}^n S_b(s)\,\Phi(t\mid s).
\] We call such a padding rule \textbf{kernel-feasible}. Mixture
compilers are a (strict) subclass that enforce an output form
\(P_b=(1-\lambda)Q+\lambda S_b\) for a single secret-independent
distribution \(Q\) (chosen to permit a coupling with \(T\ge S\) almost
surely) and a parameter \(\lambda\in[0,1]\). \textbf{Proposition 2.2
(explicit separation).} There exist success-time laws \(S_0,S_1\) and a
StopTime-IND level \(\varepsilon\) for which a kernel-feasible padding
achieves StopTime-IND\(_\varepsilon\) with strictly smaller expected
stop time than any mixture compiler. \emph{Example and proof.} Take
\(n=6\) and \[
S_0=(0.706,0.033,0,0.028,0.220,0.013),\qquad
S_1=(0.006,0.474,0.026,0.120,0.342,0.032)
\] over \(\{1,\dots,6\}\). One checks that \(\mathrm{TV}(S_0,S_1)=0.7\),
\(\mathbb{E}[S_0]=2.062\), \(\mathbb{E}[S_1]=3.414\), and that \(S_1\)
stochastically dominates \(S_0\) (tail-by-tail). Fix
\(\varepsilon=0.1\), i.e.~target
\(\mathrm{TV}(P_0,P_1)\le 2\varepsilon=0.2\). (a)
\textbf{Kernel-feasible padding.} Let \(\Phi\) be the identity for
\(s\in\{2,\dots,6\}\), and for \(s=1\) set \[
\Phi(1\mid 1)=\tfrac{2}{7},\qquad
\Phi(2\mid 1)=\tfrac{441}{700},\qquad
\Phi(3\mid 1)=\tfrac{26}{700},\qquad
\Phi(4\mid 1)=\tfrac{33}{700}.
\] Then \(T\ge S\) always. The induced laws are \[
P_0\approx(0.201714,0.477780,0.026223,0.061283,0.220000,0.013000),
\] \[
P_1\approx(0.001714,0.477780,0.026223,0.120283,0.342000,0.032000),
\] which satisfy \(\mathrm{TV}(P_0,P_1)=0.2\) (hence
StopTime-IND\(_{0.1}\)). The average expected stop time is \[
\tfrac12(\mathbb{E}[T\mid 0]+\mathbb{E}[T\mid 1])\approx 3.039074.
\] (b) \textbf{Mixture lower bound.} For any mixture compiler,
\(P_b=(1-\lambda)Q+\lambda S_b\), so \[
\mathrm{TV}(P_0,P_1)=\lambda\,\mathrm{TV}(S_0,S_1)=0.7\lambda.
\] Thus \(\lambda\le 2/7\) is necessary to meet
\(\mathrm{TV}(P_0,P_1)\le 0.2\). Moreover, since \(S_1\) stochastically
dominates \(S_0\), any \(Q\) that can be coupled to satisfy \(T\ge S\)
in both worlds must stochastically dominate \(S_1\), giving
\(\mathbb{E}[Q]\ge \mathbb{E}[S_1]=3.414\). Therefore the average
expected stop time obeys \[
\tfrac12(\mathbb{E}[T\mid 0]+\mathbb{E}[T\mid 1])
=(1-\lambda)\mathbb{E}[Q]+\lambda\tfrac{\mathbb{E}[S_0]+\mathbb{E}[S_1]}{2}
\ge (1-\lambda)\mathbb{E}[S_1]+\lambda\tfrac{\mathbb{E}[S_0]+\mathbb{E}[S_1]}{2}.
\] The right-hand side is minimized by the largest feasible
\(\lambda=2/7\) and equals \(\approx 3.220857\). Hence every mixture
compiler incurs average expected stop time at least \(3.220857\),
strictly larger than the kernel-feasible padding above. This separation
shows that letting the padding kernel depend on the realized success
time can strictly reduce overhead at a fixed StopTime-IND level. In
particular, ``mixture optimality'' does not extend beyond the mixture
class.

\hypertarget{optimal-kernel-feasible-padding-is-a-linear-program-new}{%
\subsection{2.5 Optimal kernel-feasible padding is a linear program
(new)}\label{optimal-kernel-feasible-padding-is-a-linear-program-new}}

Section 2.4 exhibited a hand-designed kernel that beats every mixture
compiler at the same StopTime-IND level. This is not an accident: in
discrete time, optimizing kernel-feasible padding under a TV constraint
is a \textbf{linear program}. \textbf{Problem 2.3 (Kernel-TV-Pad).} Fix
a horizon \(\{1,\dots,n\}\), success-time laws \(S_0,S_1\) on
\(\{1,\dots,n\}\), and a TV budget \(\tau\in[0,1]\). Choose a
kernel-feasible padding matrix \(\Phi(t\mid s)\) (Definition 2.1) to
minimize the \emph{average} expected padded stop time \[
\min_{\Phi}\ \frac12\bigl(\mathbb{E}[T\mid 0]+\mathbb{E}[T\mid 1]\bigr)
\quad\text{s.t.}\quad
\mathrm{TV}(P_0,P_1)\le \tau,
\] where \(P_b = S_b\Phi\) are the induced stop-time laws. The
constraint \(\mathrm{TV}(P_0,P_1)\le\tau\) can be linearized by
introducing auxiliary variables \(d_t\ge 0\) such that \[
 d_t \ge P_0(t)-P_1(t),\qquad d_t \ge -(P_0(t)-P_1(t)),\qquad \sum_{t=1}^n d_t \le 2\tau.
\] Since \(P_b(t)=\sum_{s\le t} S_b(s)\Phi(t\mid s)\) is linear in
\(\Phi\), all constraints and the objective are linear.
\textbf{Proposition 2.4 (LP characterization and sparsity).}
Kernel-TV-Pad is a linear program in the variables
\(\{\Phi(t\mid s):1\le s\le t\le n\}\) and \(\{d_t\}_{t=1}^n\).
Therefore an optimal solution exists at an extreme point (basic feasible
solution). In particular, there is an optimal solution with at most
\(O(n)\) nonzero variables; one crude bound from the standard ``\#basic
variables ≤ \#constraints'' principle is that some optimum uses at most
\(3n+1\) positive variables across \((\Phi,d)\). \emph{Proof sketch.}
After linearization, the feasible region is a bounded polytope defined
by \(n\) row-sum equalities, \(2n\) linear inequalities enforcing
\(|P_0(t)-P_1(t)|\le d_t\), one TV-budget inequality, and nonnegativity.
Linear objectives over polytopes attain minima at extreme points, and
basic feasible solutions have support bounded by the number of active
constraints. \textbf{Proposition 2.5 (Kernel-TV-Pad dual and KKT
interpretation).} Let \(\delta_s:= S_0(s)-S_1(s)\) and
\(w_s:=\tfrac12(S_0(s)+S_1(s))\). Consider the linearized Kernel-TV-Pad
LP from Problem 2.3. A standard dual for the \(\le\)-formulation
(Problem 2.3 with \(A_{ub}x\le b_{ub}\) and \(x\ge 0\)) uses variables:
- \(y_s\in\mathbb{R}\) for the row-sum equalities
\(\sum_{t\ge s}\Phi(t\mid s)=1\), - \(\alpha_t,\beta_t\ge 0\) for the
two inequalities \((P_0(t)-P_1(t))\le d_t\) and
\(-(P_0(t)-P_1(t))\le d_t\), - \(\gamma\ge 0\) for the TV-budget
constraint \(\sum_t d_t\le 2\tau\), and can be written as: \[
\max\ \sum_{s=1}^n y_s - 2\tau\,\gamma
\quad\text{s.t.}\quad
y_s - \delta_s\,(\alpha_t-\beta_t) \le w_s\,t\ \ \forall\ 1\le s\le t\le n,\qquad
\alpha_t+\beta_t\le \gamma\ \ \forall t.
\] Moreover, for any primal-dual optimal pair, complementary slackness
implies: - if \(\Phi(t\mid s)>0\) then
\(y_s-\delta_s(\alpha_t-\beta_t)=w_s t\); - if \(d_t>0\) then
\(\alpha_t+\beta_t=\gamma\); - if \(\gamma>0\) then the TV constraint is
tight: \(\sum_t d_t=2\tau\). \emph{Interpretation.} \(\gamma\) is the
shadow price of tightening the TV budget, and \(q_t:=\alpha_t-\beta_t\)
acts as a signed ``price schedule'' over output times \(t\).
Complementary slackness explains why optimal kernels are typically
sparse and randomize only on times \(t\) where the reduced cost is zero.
\textbf{Certificate note.} Since Kernel-TV-Pad is an LP, a dual feasible
solution matching the primal objective value certifies optimality. For
the \(n=6\) separation in Proposition 2.2, our solver returns a
primal-dual pair with primal-dual gap \(<10^{-15}\), giving a numerical
certificate of optimality for that instance. \textbf{Computational
note.} Kernel-TV-Pad can be solved efficiently for moderate \(n\) by
off-the-shelf LP solvers; in our experiments, optimal kernels are
typically sparse and randomize on only a few success-time values. We
provide a reproducibility script for the LP and a separate script that
extracts and checks dual certificates (including complementary
slackness).

\hypertarget{a-primal-dual-certificate-for-the-n6-separation-new}{%
\subsubsection{\texorpdfstring{2.5.1 A primal-dual certificate for the
\(n=6\) separation
(new)}{2.5.1 A primal-dual certificate for the n=6 separation (new)}}\label{a-primal-dual-certificate-for-the-n6-separation-new}}

For the explicit \(n=6\) instance in Proposition 2.2 with \(\tau=0.2\),
the LP optimum is attained by the kernel \[
\Phi(\cdot\mid 1)=(2/7,\ 441/700,\ 26/700,\ 33/700,\ 0,\ 0),\qquad \Phi(t\mid s)=\mathbf{1}\{t=s\}\ \text{for }s\ge 2,
\] which induces
\(P_0=(0.201714286,0.477780000,0.026222857,0.061282857,0.220000000,0.013000000)\)
and
\(P_1=(0.001714286,0.477780000,0.026222857,0.120282857,0.342000000,0.032000000)\),
with \(\mathrm{TV}(P_0,P_1)=0.2\) and objective
\(\tfrac12(\mathbb{E}[T\mid 0]+\mathbb{E}[T\mid 1])\approx 3.039074286\).
A dual certificate with \(\gamma=267/350\),
\[\alpha=(267/350,\ 89/175,\ 0,\ 0,\ 0,\ 0),\qquad
\beta=(0,\ 89/350,\ 89/350,\ 267/350,\ 267/350,\ 267/350),\] (and a
feasible \(y\in\mathbb{R}^6\) reported in the script output) achieves
the same dual objective value to machine precision. The script
\texttt{stop\_time\_dual\_cert\_v16.py} checks dual feasibility and
complementary slackness (max violation \(\approx 10^{-16}\)) and prints
both the primal kernel and dual variables.

\hypertarget{separation-families-and-empirical-scaling-new}{%
\subsection{2.6 Separation families and empirical scaling
(new)}\label{separation-families-and-empirical-scaling-new}}

Section 2.4 gives a fully explicit \(n=6\) separation. To test whether
this is a corner case, we solved Kernel-TV-Pad and the best feasible
mixture baseline across random instances (Dirichlet draws) using an LP
solver. \textbf{Proposition 2.6 (a larger explicit instance, found by LP
search).} For \(n=10\) and \(\tau=0.2\), the following pair of
success-time laws achieves a gap of about \(0.4635\) in \emph{average
expected stop time} between the best mixture compiler and the optimal
kernel-feasible padding: \[
\begin{aligned}
S_0&\approx (0.147285,0.057829,0.363867,0.036741,0.155830,0.044686,0.008706,0.030916,0.063848,0.090290),\\
S_1&\approx (0.003472,0.002795,0.000898,0.414900,0.023744,0.093175,0.035323,0.086139,0.078840,0.260714).
\end{aligned}
\] In this instance the LP optimum has average \(\approx 5.8119\), the
mixture optimum has average \(\approx 6.2754\), and
\(\mathrm{TV}(P_0,P_1)=0.2\). \textbf{Proposition 2.7 (gaps can be
\(>0.6\) at moderate horizons).} For \(n=15\) and \(\tau=0.2\), an LP
search over 400 random trials found an instance with gap
\(\approx 0.6861\) between the mixture optimum and the kernel-feasible
optimum, while still satisfying \(\mathrm{TV}(P_0,P_1)=0.2\).
\emph{Status.} These larger instances are obtained by an automated
search and are \textbf{machine-certified optimal} (primal--dual KKT
checks) by the scripts listed in Appendix C. They illustrate that
mixture suboptimality persists beyond the smallest toy example and can
grow with horizon; Section 2.7 then gives an explicit parametric family
with closed-form optima and a linear-in-horizon gap. We provide such a
family in Section 2.7.

\hypertarget{optimal-kernels-can-randomize-on-multiple-success-time-rows}{%
\subsubsection{2.6.1 Optimal kernels can randomize on multiple
success-time
rows}\label{optimal-kernels-can-randomize-on-multiple-success-time-rows}}

The closed-form family in Section 2.7 exhibits optimal kernels that
randomize only on the earliest success time. However, Kernel-TV-Pad can
produce optima that randomize on \textbf{multiple} success-time values.
This matters for implementers: an optimal padding rule may need to treat
several success-time regimes differently. \textbf{Proposition 2.7.1 (a
two-row randomization example).} Let \(n=6\), \(\tau=0.2\), and \[
S_0=(0.2795072756,\ 0.0596474087,\ 0.5803830930,\ 0.0743372489,\ 0.0059209143,\ 0.0002040595),
\] \[
S_1=(0.4861652867,\ 0.3024844985,\ 0.0949222891,\ 0.1079791870,\ 0.0066330131,\ 0.0018157256).
\] The Kernel-TV-Pad optimum induces stop-time laws \(P_0,P_1\)
satisfying \(\mathrm{TV}(P_0,P_1)=\tau\) and is attained by a kernel
that randomizes on \textbf{two} success-time rows (\(s=1\) and \(s=3\)),
e.g. \[
\Phi(\cdot\mid 1)=(0.9677824683,0,0.0322175317,0,0,0),
\] \[
\Phi(\cdot\mid 3)=(0,0,0.9259143009,0.0692989791,0.0014668514,0.0033198687),
\] with all other rows deterministic
(\(\Phi(t\mid s)=\mathbf{1}\{t=s\}\) for \(s\notin\{1,3\}\)). The
resulting average expected padded stop time is
\(\tfrac12(\mathbb{E}[T\mid 0]+\mathbb{E}[T\mid 1])\approx 2.39349\).
\emph{Verification.} The script
\texttt{stop\_time\_multirow\_cert\_v47.py} solves Kernel-TV-Pad for
this instance and verifies KKT conditions using HiGHS dual multipliers
(primal-dual gap and complementary slackness residuals
\(\lesssim 10^{-15}\)).

\hypertarget{a-closed-form-parametric-separation-family-new}{%
\subsection{2.7 A closed-form parametric separation family
(new)}\label{a-closed-form-parametric-separation-family-new}}

We next give an explicit family \(S_0,S_1\) for which (i) the
kernel-feasible optimum is attained by a simple closed-form padding
kernel, (ii) the optimal mixture compiler can be computed in closed
form, and (iii) the \textbf{mixture-vs-kernel gap grows linearly in the
horizon}. Fix a horizon \(n\ge 3\), a TV budget \(\tau\in(0,1)\), and a
parameter \(p\in(\tau,1]\). Consider the success-time laws on
\([n]=\{1,\dots,n\}\) \[
S_0:= p\,\delta_1 + (1-p)\,\delta_n,\qquad
S_1:= (p-\tau)\,\delta_2 + (1-p+\tau)\,\delta_n.
\] Then \(\mathrm{TV}(S_0,S_1)=p\). Moreover \(S_1\) first-order
stochastically dominates \(S_0\) (it shifts mass later). \textbf{Theorem
2.8 (closed-form separation; linear gap).} For the pair \(S_0,S_1\)
above and TV budget \(\tau\): 1. (\textbf{Kernel-feasible optimum}) The
optimal solution of Kernel-TV-Pad keeps all rows \(s\ge 2\) as identity
and randomizes only on row \(s=1\): \[
 \Phi(1\mid 1)=\tau/p,\qquad \Phi(2\mid 1)=1-\tau/p,
 \] with \(\Phi(t\mid s)=\mathbf{1}\{t=s\}\) for all other \((s,t)\).
The induced stop-time laws are \[
 P_0=\tau\,\delta_1+(p-\tau)\,\delta_2+(1-p)\,\delta_n,\qquad
 P_1=(p-\tau)\,\delta_2+(1-p+\tau)\,\delta_n,
 \] which satisfy \(\mathrm{TV}(P_0,P_1)=\tau\). The resulting
\textbf{average expected stop time} equals \[
 \mathrm{OPT}_{\mathrm{kern}} = \mathbb{E}[S_1] - \frac{(n-1)\tau}{2}
 = 2(p-\tau)+\frac{\tau}{2} + n\Bigl(1-p+\frac{\tau}{2}\Bigr).
 \] 2. (\textbf{Mixture-compiler optimum}) Any mixture compiler
\(C_{\lambda,Q}\) yields
\(\mathrm{TV}(P_0,P_1)=\lambda\,\mathrm{TV}(S_0,S_1)=\lambda p\), hence
feasibility forces \(\lambda\le \tau/p\). Since \(S_1\) stochastically
dominates \(S_0\), any feasible \(Q\) must stochastically dominate
\(S_1\), so \(\mathbb{E}[Q]\ge \mathbb{E}[S_1]\) with equality at
\(Q=S_1\). The mixture optimum takes \(\lambda=\tau/p\) and \(Q=S_1\),
giving \[
 \mathrm{OPT}_{\mathrm{mix}} = \mathbb{E}[S_1] - \frac{\tau}{2} - \frac{(n-2)\tau^2}{2p}.
 \] 3. (\textbf{Gap}) The suboptimality of the mixture class is
explicit: \[
 \mathrm{OPT}_{\mathrm{mix}}-\mathrm{OPT}_{\mathrm{kern}}
 = \frac{(n-2)\tau}{2}\Bigl(1-\frac{\tau}{p}\Bigr)
 \;>\;0.
 \] \emph{Proof sketch.} (1) Since \(P_1(1)=0\), TV-feasibility implies
\(P_0(1)\le \tau\). But \(P_0(1)=S_0(1)\Phi(1\mid 1)=p\Phi(1\mid 1)\),
so any feasible kernel must satisfy \(\Phi(1\mid 1)\le \tau/p\). Moving
mass in row 1 to earlier output times decreases the objective, so
optimality forces \(\Phi(1\mid 1)=\tau/p\). The remaining mass on row 1
is best placed at time 2 (the earliest time that also matches
\(P_1(2)=p-\tau\)), yielding the claimed \(\{1,2\}\)-supported row. (2)
and (3) follow from the mixture TV scaling and the fact that stochastic
dominance forces \(Q\succeq S_1\). \(\square\) \textbf{Example.} Take
\(n=6\), \(\tau=0.2\), and \(p=0.9\). Then
\(\mathrm{OPT}_{\mathrm{mix}}-\mathrm{OPT}_{\mathrm{kern}}\approx 0.3111\),
matching the LP optimum. \textbf{Proposition 2.9 (explicit dual witness
for Theorem 2.8).} In the dual of Proposition 2.5, write
\(q_t:=\alpha_t-\beta_t\) so that \(|q_t|\le \gamma\) is enforced by
\(\alpha_t+\beta_t\le \gamma\) and \(\alpha_t,\beta_t\ge 0\). For the
family in Theorem 2.8, the following choices are dual-feasible and
attain the primal value: \[
\gamma:=\tfrac14,\qquad q_1:=\tfrac14,\qquad q_t:=-\tfrac14\ \text{for all }t\in\{2,\dots,n\}.
\] Let \(y\in\mathbb{R}^n\) be \[
 y_1:=\tfrac{3p}{4},\qquad y_2:=\tfrac{5(p-\tau)}{4},\qquad y_n:=n\Bigl(1-p+\tfrac{\tau}{2}\Bigr)+\tfrac{\tau}{4},\qquad y_s:=0\ \text{for }s\notin\{1,2,n\}.
\] Setting \(\alpha_t:=\max\{q_t,0\}\) and \(\beta_t:=\max\{-q_t,0\}\)
yields \(\alpha_t,\beta_t\ge 0\) and
\(\alpha_t+\beta_t=|q_t|\le \gamma\). One can verify that the dual
constraints \(y_s-\delta_s q_t\le w_s t\) hold for every
\(1\le s\le t\le n\), with equality on the pairs used by the primal
kernel of Theorem 2.8 (namely \((1,1),(1,2),(2,2),(n,n)\)). Finally, the
dual objective equals the primal value: \[
\sum_s y_s -2\tau\gamma
= \mathrm{OPT}_{\mathrm{kern}}.
\] By strong duality, this certifies global optimality of the
closed-form kernel. A complete mechanical check of all dual inequalities
is performed by the script \texttt{stop\_time\_dual\_cert\_v18.py}.

\textbf{Figure 2 (linear gap illustration).} For the closed-form family
in Theorem 2.8, the mixture suboptimality gap is \[
\mathrm{OPT}_{\mathrm{mix}}-\mathrm{OPT}_{\mathrm{kern}}
= \frac{(n-2)\tau}{2}\Bigl(1-\frac{\tau}{p}\Bigr),
\] which grows linearly in the horizon \(n\) for fixed \(\tau,p\). The
plot below shows the gap for \(\tau=0.2\), \(p=0.9\) as a function of
\(n\).

\begin{center}
\fbox{\begin{minipage}{0.93\linewidth}\small
\textbf{Text-only replacement for Fig.~2.} The plot is omitted (text-only submission). The separation family in Theorem~2.8 gives
\[
\mathrm{OPT}_{mix}-\mathrm{OPT}_{kern}=\frac{(n-2)\tau}{2}\left(1-\frac{\tau}{p}\right),
\]
so the gap is linear in $n$ for fixed $(p,\tau)$. (To regenerate the original plot, run \texttt{\detokenize{make_figures_v67.py}}.)
\end{minipage}}
\end{center}

\hypertarget{support-complexity-beyond-single-row-randomization-new}{%
\subsection{2.8 Support complexity beyond single-row randomization
(new)}\label{support-complexity-beyond-single-row-randomization-new}}

Theorem 2.8 shows that a single randomized row can already separate
mixtures from optimal kernels. Empirically, however, optimal kernels can
randomize in many rows. \textbf{Example 2.10 (two randomized rows,
n=8).} For \(n=8\) and \(\tau=0.2\), the LP optimum for \[
\begin{aligned}
S_0 &\approx (0.339996,0.069983,0.000053,0.019840,0.055718,0.071428,0.001461,0.441521),\\
S_1 &\approx (0.005131,0.006532,0.176472,0.026986,0.095191,0.096904,0.563583,0.029202)
\end{aligned}
\] randomizes in exactly two rows, \(s\in\{1,7\}\). In particular, the
optimal row 7 splits between output times 7 and 8. The extracted dual
variables satisfy the dual constraints and match the primal objective to
numerical precision (primal-dual gap \(<10^{-15}\)). \textbf{Example
2.11 (many randomized rows, n=20).} For \(n=20\) and \(\tau=0.2\), an LP
search over random Dirichlet instances found a case where the optimal
kernel randomizes in 10 different rows while still achieving a mixture
gap of about 0.60 in average expected stop time. This suggests that
while Kernel-TV-Pad admits sparse optima (Proposition 2.4), the number
of randomized rows can still grow with \(n\). A structural
characterization of when many rows must randomize is an open problem.

\hypertarget{multi-profile-stop-time-padding-pairwise-vs-star-constraints-new}{%
\subsection{2.9 Multi-profile stop-time padding: pairwise vs star
constraints
(new)}\label{multi-profile-stop-time-padding-pairwise-vs-star-constraints-new}}

So far we presented stop-time padding for two worlds (a binary secret).
For Option B, it is useful to generalize to a \(K\)-ary secret
\(C\in[K]\) with completion-time laws \(S_c\) and padded stop-time laws
\(P_c=\Phi^\top S_c\) for a kernel-feasible padding kernel
\(\Phi(t\mid s)\). A natural \emph{pairwise} StopTime-IND\(_\tau\)
constraint is: \[
\mathrm{TV}(P_c,P_{c'}) \le \tau \qquad \forall\,c\neq c'.
\] Together with feasibility (\(\Phi(t\mid s)=0\) for \(t<s\),
\(\sum_t \Phi(t\mid s)=1\), \(\Phi\ge 0\)), minimizing average overhead
\(\frac1K\sum_c \mathbb{E}_{P_c}[T]\) remains a linear program by
introducing standard absolute-value slack variables for each pair
\((c,c')\). A simpler and often sufficient \emph{star} formulation
introduces a common reference distribution \(R\) and enforces \[
\mathrm{TV}(P_c,R) \le \tau/2 \qquad \forall\,c,
\] which implies the pairwise constraints by the triangle inequality.
The star LP uses only \(O(Kn)\) TV constraints rather than \(O(K^2n)\),
at the cost of potential conservatism. This star form also connects
directly to the posterior calculus in Section 7: if each \(P_c\) is
additively close to a common \(R\), then the stop time can be appended
as one more step in Theorem 7.2 (see Section 6.3). \textbf{Proposition
2.12 (star can be strictly suboptimal, even when feasible).} Consider
\(K=3\) profiles, horizon \(n=5\), and privacy budget \(\tau=0.2\). Let
success-time laws be \[
\begin{aligned}
S_0 &\approx (0.0005,0.6774,0.2853,0.0368,0),\\
S_1 &\approx (0.2474,0.1222,0.6093,0.0001,0.0210),\\
S_2 &\approx (0.2749,0.3145,0.0529,0.0674,0.2903).
\end{aligned}
\] Let the objective be the uniform-prior average expected stop time
\(\frac13\sum_{c=0}^2 \mathbb{E}_{P_c}[T]\). (a) Under the
\emph{pairwise} constraints \(\mathrm{TV}(P_c,P_{c'})\le \tau\) for all
\(c\neq c'\), the LP optimum is \(\mathrm{OPT}_{pair}\approx 3.08872\),
achieved by the kernel \[
\Phi_{pair} \approx
\begin{pmatrix}
0.039647 & 0.960353 & 0 & 0 & 0\\
0 & 0.787304 & 0.212696 & 0 & 0\\
0 & 0 & 0.432964 & 0.120956 & 0.446079\\
0 & 0 & 0 & 1 & 0\\
0 & 0 & 0 & 0 & 1
\end{pmatrix}.
\] (b) Under the \emph{star} constraints
\(\mathrm{TV}(P_c,R)\le \tau/2\) for a common reference \(R\), the LP
optimum is \(\mathrm{OPT}_{star}\approx 3.22584\), achieved by \[
\Phi_{star} \approx
\begin{pmatrix}
0 & 1 & 0 & 0 & 0\\
0 & 0.624820 & 0.262765 & 0.005617 & 0.106798\\
0 & 0 & 0.450269 & 0.122898 & 0.426833\\
0 & 0 & 0 & 1 & 0\\
0 & 0 & 0 & 0 & 1
\end{pmatrix},
\qquad
R\approx (0,0.423753,0.206459,0.075668,0.294120).
\] Therefore the star formulation is strictly conservative here:
\(\mathrm{OPT}_{star}-\mathrm{OPT}_{pair}\approx 0.13711\).
\emph{Verification.} The script \texttt{stop\_time\_star\_gap\_v21.py}
solves both LPs and prints the optimal kernels and induced TVs;
\texttt{stop\_time\_kary\_cert\_v38.py} additionally extracts HiGHS dual
multipliers and checks KKT. \textbf{Certification remark.} Since both
the pairwise and star formulations are linear programs, their optima
admit short primal-dual certificates. The script
\texttt{stop\_time\_kary\_cert\_v38.py} reruns the instance of
Proposition 2.12, extracts the HiGHS dual multipliers, and checks the
KKT conditions (primal feasibility, reduced-cost nonnegativity,
complementary slackness), reporting a primal-dual gap on the order of
\(10^{-15}\). This makes the numerical separation fully
machine-checkable. Thus, while the star formulation reduces the number
of TV constraints from \(O(K^2n)\) to \(O(Kn)\), it can be strictly
suboptimal. We leave a tight characterization of when the star
constraints are lossless (e.g., under geometric or convexity conditions
on the feasible \{P\_c\}) as an open problem. \textbf{Lemma 2.13 (TV
diameter vs radius; barycentric star bound is optimal).} Let
\(P_1,\dots,P_K\) be distributions on a finite alphabet, and let \[
D:= \max_{i,j\in[K]} \mathrm{TV}(P_i,P_j)
\] be their \textbf{diameter} in total variation. Define the
\textbf{barycenter} \(R:= \frac{1}{K}\sum_{j=1}^K P_j\). Then for every
\(i\in[K]\), \[
\mathrm{TV}(P_i,R) \le \frac{K-1}{K}\,D.
\] Moreover, the constant \((K-1)/K\) is best possible in general: for
the \(K\) point-masses \(P_i=\delta_i\) on alphabet \(\{1,\dots,K\}\),
we have \(D=1\) but \(\min_R\max_i\mathrm{TV}(P_i,R)=1-1/K=(K-1)/K\).
\emph{Proof.} By convexity of TV in its second argument, \[
\mathrm{TV}(P_i,R)
= \mathrm{TV}\Bigl(P_i,\ \frac{1}{K}P_i + \frac{1}{K}\sum_{j\ne i} P_j\Bigr)
\le \frac{1}{K}\sum_{j\ne i} \mathrm{TV}(P_i,P_j)
\le \frac{K-1}{K}\,D.
\] Tightness follows by direct calculation for the simplex vertices.
\(\square\)

\textbf{Figure 4 (TV-star relaxation across \(\tau\)).} For the fixed
\(K=3,n=5\) instance of Proposition 2.12, we sweep \(\tau\in[0,0.5]\)
and solve both LPs. The TV-star optimum remains above the pairwise-TV
optimum across an interval, illustrating \emph{persistent} conservatism
(not just a single tuned parameter point). The figure is generated by
\texttt{make\_figures\_v67.py}.

\begin{center}
\fbox{\begin{minipage}{0.93\linewidth}\small
\textbf{Text-only replacement for Fig.~4.} The plot is omitted (text-only submission). It swept $\tau\in[0,0.5]$ for the $K{=}3,n{=}5$ instance in Proposition~2.12, showing a persistent interval where the TV-star relaxation stays above the pairwise-TV optimum. (To regenerate the original plot, run \texttt{\detokenize{make_figures_v67.py}}.)
\end{minipage}}
\end{center}

\emph{Discussion (why the TV-star shortcut can cost overhead).} The
TV-star constraints \(\mathrm{TV}(P_c,R)\le \tau/2\) imply pairwise TV
\(\le \tau\) by the triangle inequality. However, for \(K\ge 3\) they
can be much stronger than the pairwise constraints: if a family has
pairwise diameter \(\tau\), then Lemma 2.13 guarantees the existence of
a common center only within radius \((K-1)\tau/K\) (not \(\tau/2\)), and
the worst-case gap factor \[
\frac{(K-1)\tau/K}{\tau/2}=2-2/K
\] approaches \(2\) as \(K\) grows. In Section 10 we therefore treat
star-style constraints as a \textbf{typing rule} enabling clean
posterior composition, and we quantify the overhead cost separately
(Proposition 2.12).

\hypertarget{posterior-friendly-stop-time-padding-l-star-constraints-and-probability-floors-new}{%
\subsection{2.10 Posterior-friendly stop-time padding: L∞-star
constraints and probability floors
(new)}\label{posterior-friendly-stop-time-padding-l-star-constraints-and-probability-floors-new}}

Section 6.3 composes the stop time with the sequential posterior
calculus (Theorem 7.2) by treating the stop time as one more observation
step. That calculus requires a \textbf{coordinate-wise} additive
deviation bound relative to a secret-independent reference with a
\textbf{positive lower bound}. This motivates a stop-time padding
variant whose constraints are tailored to posterior control.
\textbf{Problem 2.14 (PosteriorStopTime-Pad).} Fix a horizon
\(\{1,\dots,n\}\), success-time laws \(\{S_c\}_{c\in[K]}\), and
parameters \((\kappa_T,\delta_T)\) with \(0<\kappa_T<1\) and
\(0\le \delta_T<\kappa_T\). Choose a kernel-feasible padding matrix
\(\Phi\) and a reference distribution \(R\) on \(\{1,\dots,n\}\) to
minimize average overhead \[
\min_{\Phi,R}\ \frac1K\sum_{c=1}^K \mathbb{E}_{P_c}[T]
\quad\text{s.t.}\quad
P_c = S_c\Phi,\ \Phi\ \text{kernel-feasible},\ \sum_t R(t)=1,\ R(t)\ge \kappa_T\ \forall t,
\] \[
|P_c(t)-R(t)|\le \delta_T\ \ \forall c,t.
\] This is a linear program: the constraints
\(|P_c(t)-R(t)|\le \delta_T\) are two linear inequalities per \((c,t)\),
and \(R(t)\ge \kappa_T\) is linear. \textbf{Proposition 2.15
(PosteriorStopTime-Pad plugs into Theorem 7.2).} If \((\Phi,R)\) is
feasible for Problem 2.14, then the padded stop time can be appended as
an extra observation step with \(\kappa_{T+1}=\kappa_T\) and
\(\delta_{T+1}=\delta_T\) in Theorem 7.2, yielding the posterior
inflation factor
\(\log\!\big((\kappa_T+\delta_T)/(\kappa_T-\delta_T)\big)\) contributed
by the stop time. \emph{Remark (certification).} Problem 2.14 is a
linear program with a standard primal--dual optimum certificate. In
particular, solvers such as SciPy/HiGHS return inequality marginals and
equality marginals that can be checked against the KKT conditions (dual
feasibility, complementary slackness, and a small primal--dual gap). For
a small hard-coded instance, the script
\texttt{stop\_time\_posterior\_cert\_v46.py} builds
PosteriorStopTime-Pad, solves it, extracts HiGHS dual multipliers, and
reports max KKT residuals on the order of \(10^{-16}\). \emph{Remark
(relation to TV-star constraints).} If one uses a TV-star constraint
\(\mathrm{TV}(P_c,R)\le \tau/2\), then \(\|P_c-R\|_1\le \tau\) and hence
\(\|P_c-R\|_\infty\le \tau\), so Problem 2.14 holds with
\(\delta_T=\tau\) provided one also enforces a floor
\(R(t)\ge \kappa_T\). However, \(\|\cdot\|_\infty\) control can be much
tighter than what TV alone implies, and we found concrete instances
where TV-star constraints are strictly conservative relative to pairwise
TV (Proposition 2.12).

\hypertarget{mark-leakage-collapsing-shard-marks-by-public-scheduling-psc}{%
\section{3. Mark leakage: collapsing shard marks by public scheduling
(PSC)}\label{mark-leakage-collapsing-shard-marks-by-public-scheduling-psc}}

PSC fixes a public mapping from slot index \(s\) to a shard subset
\(J_s\) and requires fixed-size traffic to every shard in \(J_s\) each
slot (real or dummy). Then ``which shard(s) were contacted'' becomes a
deterministic function of \((\text{epoch},s)\). PSC collapses
shard-level marks: conditioned on the public schedule, the set of
contacted shards is a deterministic function of time, so only the
within-set choice remains to be hidden. \textbf{Remark 3.1 (support
containment and finite KL).} PSC restricts each slot to a public
scheduled action set. If the cover rule assigns strictly positive
probability to every action in the scheduled set, then every
secret-dependent per-slot law has support contained in the cover
support. This prevents the pathological case
\(D_{\mathrm{KL}}(P\|Q)=\infty\) from support mismatch and makes the KL-
and MI-based accounting in Sections 5--6 well-defined. ---

\hypertarget{equalizing-within-a-scheduled-action-set-mc-eq-gives-a-tv-knob}{%
\section{4. Equalizing within a scheduled action set: MC-EQ gives a TV
knob}\label{equalizing-within-a-scheduled-action-set-mc-eq-gives-a-tv-knob}}

\hypertarget{tv-shrinkage-under-mixture-equalization}{%
\subsection{4.1 TV shrinkage under mixture
equalization}\label{tv-shrinkage-under-mixture-equalization}}

Let \(Q\) be a cover distribution and \(R\) a client's true-target
distribution over the same scheduled action set. If one lookup selects
one action from \(R\) and \(w-1\) actions from \(Q\), then the induced
per-lookup contact distribution is \[
P = (1/w)\,R + (1-1/w)\,Q,
\] and \[
\mathrm{TV}(P,Q)=\frac{1}{w}\mathrm{TV}(R,Q).
\] \textbf{Corollary 4.1 (history-uniform TV from one-real-among-w
mixing).} Let \(K(x,\cdot)\) be a reference ``cover'' distribution (or a
Markov kernel row) that does not depend on the secret profile \(C\).
Suppose the realized per-step action distribution conditioned on an
arbitrary history \(h_{t-1}\) and secret \(C=c\) has the mixture form \[
P_{c,h_{t-1}}(\cdot) = (1-\rho)\,K(x,\cdot) + \rho\,R_{c,h_{t-1}}(\cdot)
\] for some \(\rho\in[0,1]\), some (possibly history- and
secret-dependent) distribution \(R_{c,h_{t-1}}\), and where \(x\) is the
public state used to index the cover row (often \(x\) is simply the
previous action). Then for every \(c\) and every history \(h_{t-1}\), \[
\mathrm{TV}\bigl(P_{c,h_{t-1}},\ K(x,\cdot)\bigr) = \rho\,\mathrm{TV}\bigl(R_{c,h_{t-1}},\ K(x,\cdot)\bigr) \le \rho.
\] In particular, in the common \(\rho=1/w\) ``one real among \(w\)''
regime, the \textbf{history-uniform TV deviation} parameter satisfies
\(\delta_{eq}\le 1/w\) \emph{automatically}, without any additional
assumptions on how \(R_{c,h}\) depends on the past. \emph{Proof.}
Immediate from the convexity identity
\(\mathrm{TV}((1-\rho)Q+\rho R,Q)=\rho\,\mathrm{TV}(R,Q)\le \rho\).
\textbf{Lemma 4.2 (pointwise deviation under one-real-among-\(w\)
mixing).} Let \(K\) be a reference distribution on a finite set \(J\),
let \(R\) be any distribution on \(J\), and define \[
P:= (1-\rho)K + \rho R.
\] Then for every \(y\in J\), \[
|P(y)-K(y)| = \rho\,|R(y)-K(y)| \le \rho.
\] If additionally \(K\) is uniform on \(|J|=m\), then
\(\|P-K\|_\infty \le \rho(1-1/m) < \rho\). In the common \(\rho=1/w\)
regime, this provides a per-step additive deviation parameter
\(\delta_\infty \le 1/w\) for Theorem 7.2 and Theorem 10.1.
\emph{Proof.} Immediate from \(P-K=\rho(R-K)\) and \(|R(y)-K(y)|\le 1\)
(or \(\le 1-1/m\) in the uniform case). ---

\hypertarget{new-exact-and-moment-kl-bounds-for-one-real-among-w-injection}{%
\section{\texorpdfstring{5. New: exact and moment KL bounds for
``one-real-among-\(w\)''
injection}{5. New: exact and moment KL bounds for ``one-real-among-w'' injection}}\label{new-exact-and-moment-kl-bounds-for-one-real-among-w-injection}}

\hypertarget{exact-per-step-kl-for-injection-into-a-reference-distribution}{%
\subsection{5.1 Exact per-step KL for injection into a reference
distribution}\label{exact-per-step-kl-for-injection-into-a-reference-distribution}}

Let \(q\) be a reference distribution on a finite set \(A\) with
\(q_j>0\). Fix \(\rho=1/w\), and define \[
p_j:= (1-\rho)q + \rho\,\delta_j.
\] \textbf{Proposition 5.1 (exact KL).} For every \(j\), \[
D_{\mathrm{KL}}(p_j\|q)
= \big((1-\rho)q_j+\rho\big)\log\!\Bigl(1-\rho+\frac{\rho}{q_j}\Bigr)
+ (1-\rho)(1-q_j)\log(1-\rho).
\] \emph{Proof.} Split the KL sum into coordinate \(j\) and
\(A\setminus\{j\}\); for \(y\neq j\), \(p_j(y)=(1-\rho)q(y)\), giving
the common \(\log(1-\rho)\) term. \textbf{Uniform special case.} If
\(q\) is uniform on \(|A|=m\), then \[
D_{\mathrm{KL}}(p_j\|q)=\frac{1+\rho(m-1)}{m}\log(1+\rho(m-1))+\frac{(m-1)(1-\rho)}{m}\log(1-\rho)
=\frac{m-1}{2w^2}+O\!\left(\frac{m^2}{w^3}\right).
\]

\hypertarget{a-moment-upper-bound-avoiding-a-global-q_min}{%
\subsection{\texorpdfstring{5.2 A moment upper bound avoiding a global
\(q_{\min}\)}{5.2 A moment upper bound avoiding a global q\_\{\textbackslash min\}}}\label{a-moment-upper-bound-avoiding-a-global-q_min}}

\textbf{Lemma 5.2 (moment bound).} For \(p_j=(1-\rho)q+\rho\delta_j\),
\[
D_{\mathrm{KL}}(p_j\|q)\ \le\ \rho^2\Bigl(\frac{1}{q_j}-1\Bigr).
\] \emph{Proof.} Use \(\log(1+x)\le x\) for \(x\ge0\) and
\(\log(1-\rho)\le -\rho\); the linear-in-\(\rho\) terms cancel.
\textbf{Corollary 5.3.} If the injected index \(J\) is random, then \[
\mathbb{E}[D_{\mathrm{KL}}(p_J\|q)] \le \rho^2\left(\mathbb{E}\Bigl[\frac{1}{q_J}\Bigr]-1\right).
\] This suggests designing schedules/covers so that ``real'' contacts
land on higher-mass actions under \(q\). ---

\hypertarget{dependency-robust-mi-accounting-and-fano-inversion}{%
\section{6. Dependency-robust MI accounting and Fano
inversion}\label{dependency-robust-mi-accounting-and-fano-inversion}}

The posterior-first path (Theorem 10.1) gives
\textbf{transcript-by-transcript} guarantees but can be conservative
over long horizons. This section provides a complementary
\textbf{average-information} tool: an MI upper bound that remains valid
even when the real process depends arbitrarily on the full history, as
long as each step is TV-close to a secret-independent Markov reference
kernel.

\hypertarget{markov-reference-mi-bound}{%
\subsection{6.1 Markov-reference MI
bound}\label{markov-reference-mi-bound}}

Let \(C\) be a hidden profile and let \(O_{1:T}\) be the observed mark
process (after PSC has reduced marks to within-set choices). For each
\(c\), write \(P_c\) for the law of \(O_{1:T}\) conditioned on \(C=c\);
\(P_c\) may be non-Markov and may depend on the full history.

Fix a \textbf{reference} Markov chain \(Q\) on the mark space
\(\mathcal X\) with initial distribution \(\mu\) and kernel \(K\): \[
Q(o_{1:T})=\mu(o_1)\prod_{t=2}^T K(o_{t-1},o_t).
\] Assume a positive floor \[
k_{\min}:=\min_{x,y\in\mathcal X}K(x,y) > 0.
\] Assume the following \textbf{history-uniform TV deviation} contract
holds for the real process: for every \(t\ge 2\), every profile \(c\),
every history \(h_{t-1}\) whose last state is \(x\), and every
measurable set \(A\subseteq\mathcal X\), \[
\mathrm{TV}\bigl(P(O_t\in\cdot \mid C=c,H_{t-1}=h_{t-1}),\ K(x,\cdot)\bigr)\ \le\ \delta_{eq}.
\tag{HUTV}
\] (Under one-real-among-\(w\) MC-EQ, Corollary 4.1 gives
\(\delta_{eq}\le 1/w\) automatically.)

\textbf{Theorem 6.1 (MI via a Markov reference).} Under (HUTV) with
\(\delta_{eq}<k_{\min}\), \[
I(C;O_{1:T}) \ \le\ I(C;O_1)\ +\ (T-1)\ln\!\left(1+\frac{2\delta_{eq}^2}{k_{\min}}\right).
\] More generally, for any choice of reference initial law \(\mu\), \[
I(C;O_{1:T}) \ \le\ \mathbb E_c\!\left[D_{\mathrm{KL}}(P_c(O_1)\,\|\,\mu)\right]\ +\ (T-1)\ln\!\left(1+\frac{2\delta_{eq}^2}{k_{\min}}\right).
\]

\emph{Proof.} Fix any reference distribution \(Q\) that does not depend
on \(C\). A standard variational bound gives \[
I(C;O_{1:T}) = \mathbb E_c\!\left[D_{\mathrm{KL}}(P_c \,\|\, P)\right]\ \le\ \mathbb E_c\!\left[D_{\mathrm{KL}}(P_c \,\|\, Q)\right],
\] where \(P\) is the unconditional law of \(O_{1:T}\) and we used that
\(D_{\mathrm{KL}}(\cdot\|\cdot)\) is minimized at \(P\).

Now apply the chain rule for KL: \[
D_{\mathrm{KL}}(P_c \,\|\, Q)
= D_{\mathrm{KL}}(P_c(O_1)\,\|\,\mu) + \sum_{t=2}^T \mathbb E_{h_{t-1}\sim P_c}\!\left[D_{\mathrm{KL}}\!\left(P(O_t\mid c,h_{t-1})\,\|\,K(x,\cdot)\right)\right],
\] where \(x\) is the last state of \(h_{t-1}\).

For each conditional term, (HUTV) gives \(\mathrm{TV}\le \delta_{eq}\)
and the reference row has minimum mass \(\ge k_{\min}\), so Lemma 10.3
yields \[
D_{\mathrm{KL}}\!\left(P(O_t\mid c,h_{t-1})\,\|\,K(x,\cdot)\right)\ \le\ \ln\!\left(1+\frac{2\delta_{eq}^2}{k_{\min}}\right),
\] uniformly over \(c\) and \(h_{t-1}\). Summing over \(t=2,\dots,T\)
and taking \(\mathbb E_c\) gives the second displayed bound.

To obtain the first bound, choose \(\mu:=P(O_1)\); then
\(\mathbb E_c[D_{\mathrm{KL}}(P_c(O_1)\|P(O_1))]=I(C;O_1)\). \(\Box\)

\hypertarget{fano-inversion-and-an-explicit-w-rule}{%
\subsection{\texorpdfstring{6.2 Fano inversion and an explicit \(w\)
rule}{6.2 Fano inversion and an explicit w rule}}\label{fano-inversion-and-an-explicit-w-rule}}

If \(C\sim\mathrm{Unif}([K])\) and \(I(C;O_{1:T})\le B\), then Fano's
inequality implies \[
P_{\mathrm{succ}} := \Pr[\widehat C(O_{1:T})=C] \ \le\ \frac{B+\ln 2}{\ln K}.
\] Combining this with Theorem 6.1 and the MC-EQ plug-in
\(\delta_{eq}\le 1/w\) yields the explicit sufficient condition \[
\frac{(T-1)\ln\!\left(1+\frac{2}{k_{\min}w^2}\right)+I(C;O_1)+\ln 2}{\ln K}\ \le\ \alpha
\quad\Rightarrow\quad
P_{\mathrm{succ}}\le \alpha,
\] which can be inverted for \(w\) given a target success probability
\(\alpha\). For small \(1/(k_{\min}w^2)\) this behaves like \[
w \ \gtrsim\ \sqrt{\frac{2(T-1)}{k_{\min}\,(\alpha\ln K-\ln 2-I(C;O_1))}}.
\] In contrast to posterior-first composition (Theorem 10.1), this
MI/Fano path scales roughly like \(\sqrt{T}\) in the high-\(w\) regime.
---

\hypertarget{end-to-end-compiler-calculus-folding-stop-time-into-the-same-posterior-calculus-new}{%
\subsection{6.3 End-to-end compiler calculus: folding stop time into the
same posterior calculus
(new)}\label{end-to-end-compiler-calculus-folding-stop-time-into-the-same-posterior-calculus-new}}

Option B's goal is to reason about \emph{joint} transcripts that include
both (i) a contact-mark trace and (ii) a termination time. A convenient
way to make these modules compose is to treat the stop time as just
\textbf{one more observation step} and apply the sequential posterior
bound from Theorem 7.2. Let \(C\sim \mathrm{Unif}([K])\). Consider an
observation sequence \[
(O_1,\dots,O_T,\; \widetilde T),
\] where \(O_{1:T}\) is a contact-mark trace and \(\widetilde T\) is the
padded stop time produced by a stop-time compiler. Assume: 1.
(\textbf{Mark compiler hypothesis}) For each step \(t\le T\) there is a
secret-independent reference \(K_t(\cdot\mid h)\) with
\(\kappa_t:= \inf_{h,y} K_t(y\mid h) > 0\) and an additive deviation
bound \[
\bigl|P(O_t=y\mid C=c,H_{t-1}=h) - K_t(y\mid h)\bigr| \le \delta_t < \kappa_t.
\] 2. (\textbf{Stop-time compiler hypothesis}) There is a
secret-independent reference \(Q\) on stop times with
\(\kappa_{T+1}:=\min_y Q(y) > 0\) such that for all \(c\) and all
histories \(h\), \[
\bigl|P(\widetilde T=y\mid C=c, H_T=h) - Q(y)\bigr| \le \delta_{T+1} < \kappa_{T+1}.
\] Then Theorem 7.2 applies to the length-\(T{+}1\) transcript and
yields, for every realized transcript, \[
\max_c P\bigl(C=c\mid O_{1:T},\widetilde T\bigr) \le \exp(\varepsilon_{T+1})/K,
\quad
\varepsilon_{T+1}=\sum_{t=1}^{T+1}\log\!\left(\frac{\kappa_t+\delta_t}{\kappa_t-\delta_t}\right).
\] \textbf{Concrete plug-ins.} - Under PSC with uniform scheduled sets
\(|J_t|=m_t\), the one-real-among-\(w\) mixing form together with Lemma
4.2 gives the pointwise deviation bound \(\delta_t\le 1/w\).
Consequently \(\varepsilon_T \approx \sum_t 2m_t/w\) when \(m_t/w\) is
small. - Under a \emph{mixture} stop-time compiler that outputs
\(\widetilde T\sim Q\) with probability \(1-\lambda\) and otherwise
outputs an arbitrary \(\widetilde T\) (subject to feasibility), one has
the additive deviation bound \(\delta_{T+1}\le \lambda\) for every \(y\)
because \(P(\cdot)-Q(\cdot)=\lambda(R(\cdot)-Q(\cdot))\) with
\(|R(y)-Q(y)|\le 1\). This yields an explicit end-to-end posterior
control inequality coupling the two knobs \(w\) (mark equalization) and
\(\lambda\) (stop-time mixing). In particular, a sufficient condition
for \(\max_c P(C=c\mid \text{transcript}) \le (1+\eta)/K\) is
\(\varepsilon_{T+1} \le \log(1+\eta)\), i.e.~a sum of one term from each
step.

\hypertarget{new-near-1k-identification-via-pointwise-maximal-leakage-from-tv-row-perturbation}{%
\section{\texorpdfstring{7. New: near-\(1/K\) identification via
pointwise maximal leakage from TV-row
perturbation}{7. New: near-1/K identification via pointwise maximal leakage from TV-row perturbation}}\label{new-near-1k-identification-via-pointwise-maximal-leakage-from-tv-row-perturbation}}

\hypertarget{pml-definition-and-a-posterior-bound-under-a-uniform-prior}{%
\subsection{7.1 PML definition and a posterior bound under a uniform
prior}\label{pml-definition-and-a-posterior-bound-under-a-uniform-prior}}

For a channel \(W_u(y\mid c)\) (at fixed state \(u\)), define the
pointwise maximal leakage under a uniform prior \(P_C\) as \[
\mathrm{PML}_u:= \sup_y \log\frac{\max_c W_u(y\mid c)}{\sum_c P_C(c)\,W_u(y\mid c)}.
\] This is a pointwise (per-output) variant of \textbf{maximal leakage}
{[}IWM18{]}. The following identity makes it directly interpretable as a
posterior bound: Under a uniform prior,
\(\max_c \Pr[C=c\mid y] = e^{\mathrm{PML}_u}/K\). Hence bounding
\(\mathrm{PML}_u\) gives a sharp near-\(1/K\) identification guarantee.

\hypertarget{additive-perturbation-finite-pml-and-thus-near-1k-posterior-control}{%
\subsection{\texorpdfstring{7.2 Additive perturbation ⇒ finite PML (and
thus near-\(1/K\) posterior
control)}{7.2 Additive perturbation ⇒ finite PML (and thus near-1/K posterior control)}}\label{additive-perturbation-finite-pml-and-thus-near-1k-posterior-control}}

Assume a reference row \(K(u,\cdot)\) has a uniform lower bound
\(\kappa:= \min_y K(u,y)\), and assume an additive deviation bound
\(\delta_\infty:= \sup_{c,y}|W_u(y\mid c)-K(u,y)|\). If
\(\kappa>2\delta_\infty\), then \[
\mathrm{PML}_u \le \sup_y \log\frac{K(u,y)+\delta_\infty}{K(u,y)-\delta_\infty}\le \log\frac{\kappa+\delta_\infty}{\kappa-\delta_\infty}.
\] Moreover, in general
\(\delta_\infty\le 2\,\mathrm{TV}(W_u(\cdot\mid c),K(u,\cdot))\). Under
the concrete one-real-among-\(w\) mixing form
\(W_u=(1-1/w)K(u,\cdot)+(1/w)R\), Lemma 4.2 sharpens this to
\(\delta_\infty\le 1/w\). Thus for \(\kappa=\Theta(1/d)\) (lazy
bounded-degree cover), \(\mathrm{PML}_u\approx O(d/w)\).
\textbf{Corollary 7.1 (explicit \(w\) for near-\(1/K\) posterior).} Fix
\(\eta\in(0,1)\) and target \(\max_c \Pr[C=c\mid y]\le (1+\eta)/K\)
(i.e., within a factor \(1+\eta\) of random guessing). It suffices that
\(\mathrm{PML}_u\le \ln(1+\eta)\), i.e. \[
\frac{\kappa+\delta_\infty}{\kappa-\delta_\infty}\ \le\ 1+\eta
\quad\Longleftrightarrow\quad
\delta_\infty \le \frac{\eta}{2+\eta}\,\kappa.
\] Under the plug-in \(\delta_\infty\le 1/w\) (Lemma 4.2), a sufficient
condition is \[
w \ \ge\ \frac{(2+\eta)}{\eta\,\kappa}.
\] In particular, for a lazy bounded-degree cover kernel with
\(\kappa\ge 1/(2d)\), it suffices that \[
w \ \ge\ \frac{2d(2+\eta)}{\eta}\ \approx\ \frac{4d}{\eta}.
\] \emph{Remark.} This route is designed for the regime where one wants
guarantees close to random guessing \(1/K\); MI→Fano can be vacuous
there. ---

\hypertarget{new-sequential-posterior-composition-under-per-history-additive-deviation}{%
\subsection{7.3 New: sequential posterior composition under per-history
additive
deviation}\label{new-sequential-posterior-composition-under-per-history-additive-deviation}}

The per-step PML bound above is local (one step / one state). For a full
transcript \(y_{1:T}\) we can control the posterior directly under a
mild \textbf{per-history additive deviation} hypothesis. \textbf{Theorem
7.2 (Sequential posterior bound from per-history additive deviation).}
Assume \(C\) is uniform on \([K]\). For each time \(t\) and each history
\(h_{t-1}\), let \(K_t(\cdot\mid h_{t-1})\) be a reference distribution
that does not depend on \(C\), and define -
\(kappa_t:= \inf_{h,y} K_t(y\mid h)\). Suppose the real conditional laws
satisfy, for all secrets \(c\), histories \(h\), and outcomes \(y\), -
\(|P(O_t=y\mid C=c, H_{t-1}=h) - K_t(y\mid h)| \le delta_t\), with
\(0 \le delta_t < kappa_t\). Then for every transcript \(y_{1:T}\), -
\(\max_c P(C=c\mid O_{1:T}=y_{1:T}) \le \exp(epsilon_T)/K\), where -
\(epsilon_T:= \sum_{t=1}^T \log((kappa_t+delta_t)/(kappa_t-delta_t))\).
\emph{Proof.} For any \(c,c'\) and any history \(h\), the hypothesis
implies \(P(O_t=y\mid c,h) \le K_t(y\mid h)+delta_t\) and
\(P(O_t=y\mid c',h) \ge K_t(y\mid h)-delta_t\), so the conditional
likelihood ratio is bounded by \((kappa_t+delta_t)/(kappa_t-delta_t)\).
Multiply these bounds over time (chain rule for likelihoods) to bound
the transcript likelihood ratio. With a uniform prior, Bayes' rule
implies the stated posterior maximum. \textbf{Corollary 7.3 (PSC + MC-EQ
plug-in; linear \(w\) budget).} If \(K_t\) is uniform on a scheduled set
\(J_t\) of size \(m_t\), then \(kappa_t = 1/m_t\). If one only knows a
per-history TV deviation bound
\(\,\mathrm{TV}(P(\cdot\mid c,h),K_t(\cdot\mid h))\le \delta_{eq,t}\),
then the pointwise deviation satisfies the generic inequality \[
\delta_t:= \sup_{c,h,y}\bigl|P(O_t=y\mid c,h)-K_t(y\mid h)\bigr| \le 2\delta_{eq,t}.
\] In the common one-real-among-\(w\) mixing form on \(J_t\), \[
P(\cdot\mid c,h)=(1-1/w)K_t(\cdot\mid h) + (1/w)R_{c,h}(\cdot),
\] Lemma 4.2 yields the sharper bound \(\delta_t\le 1/w\). Provided
\(w>m_t\) for all \(t\), Theorem 7.2 gives \[
\epsilon_T \le \sum_{t=1}^T \log\!\left(\frac{1/m_t+1/w}{1/m_t-1/w}\right).
\] For small \(m_t/w\), this is approximately \(\sum_t 2m_t/w\).
Therefore, to ensure \(\max_c P(C=c\mid y) \le (1+eta)/K\), it suffices
to have \(\epsilon_T \le \log(1+eta)\), giving the rule \[
w \gtrsim (2/\log(1+eta)) \sum_{t=1}^T m_t \approx (2/eta) \sum_{t=1}^T m_t.
\] \emph{Remark.} This is a worst-case composition bound. In
long-horizon regimes, MI/KL accounting can scale like \(\sqrt{T}\) for
the same primitive, while Theorem 7.2 scales linearly in the number of
perturbed steps. Theorem 7.2 is most useful for short sessions, or when
one can bound the number of ``real injections'' explicitly.

\hypertarget{binary-distinguishing-robust-kltv-envelopes-and-non-vacuous-design-rules}{%
\section{8. Binary distinguishing: robust KL→TV envelopes and
non-vacuous design
rules}\label{binary-distinguishing-robust-kltv-envelopes-and-non-vacuous-design-rules}}

To upper bound \(\mathrm{TV}(P,Q)\) from a KL bound
\(D:=D_{\mathrm{KL}}(P\|Q)\), record Pinsker and Bretagnolle-Huber and
the robust envelope \[
g(D):= \min\{\sqrt{D/2},\ \sqrt{1-e^{-D}}\},
\qquad \mathrm{TV}(P,Q)\le g(D).
\] Consequently, any binary distinguisher between two transcript laws
\(P,Q\) with \(D_{\mathrm{KL}}(P\|Q)\le D\) has advantage at most
\(g(D)\). In particular, to ensure \(\mathrm{TV}(P,Q)\le \varepsilon\),
it suffices that \[
D \le 2\varepsilon^2 \quad\text{(Pinsker)}\qquad\text{or}\qquad D \le -\ln(1-\varepsilon^2) \quad\text{(Bretagnolle--Huber)}.
\] The Bretagnolle--Huber inversion remains informative for moderate
\(\varepsilon\) where the small-\(D\) Pinsker regime can be overly
strict. ---

\hypertarget{mixing-time-stability-under-equalization}{%
\section{9. Mixing-time stability under
equalization}\label{mixing-time-stability-under-equalization}}

Section 6 bounds \emph{privacy} leakage of the observation sequence
relative to a Markov reference. Here we bound a different failure mode:
equalization perturbs the intended cover process, potentially destroying
the \emph{routing/latency} benefits (mixing and stationary behavior) of
the cover Markov chain. We model the scheduled cover process as a Markov
kernel K on a finite state space X (e.g., committees). The realized
process uses a kernel K\_c (possibly secret-dependent) that is row-wise
close to K. \textbf{Assumption (row-wise TV perturbation).} Let δ\_eq:=
sup\_x TV(K\_c(x,·), K(x,·)). In MC-EQ-type mechanisms, δ\_eq typically
scales like Θ(1/w).

\hypertarget{one-step-contractive-case-ux3b21}{%
\subsection{9.1 One-step contractive case
(β\textless1)}\label{one-step-contractive-case-ux3b21}}

Define the Dobrushin (TV contraction) coefficient β(P):= sup\_\{x,x'\}
TV(P(x,·), P(x',·)). Then for any distributions μ,ν, we have
\textbar\textbar μP − νP\textbar\textbar\_TV ≤ β(P) \textbar\textbar μ −
ν\textbar\textbar{}\emph{TV. \textbf{Lemma 9.1 (geometric perturbation
recursion).} Let Δ\_t:= \textbar\textbar{} μ K\_c\^{}t − μ K\^{}t
\textbar\textbar{}\emph{TV. If β(K)=β\textless1, then Δ}\{t+1\} ≤ δ\_eq
+ β Δ\_t, hence Δ\_t ≤ δ\_eq (1−β\^{}t)/(1−β) ≤ δ\_eq/(1−β). \emph{Proof
sketch.} Add/subtract (μ K\_c\^{}t)K: Δ}\{t+1\} =
\textbar\textbar μK\_c\^{}tK\_c − μK\^{}tK\textbar\textbar{} ≤
\textbar\textbar μK\_c\^{}t(K\_c−K)\textbar\textbar{} +
\textbar\textbar(μK\_c\^{}t)K − (μK\^{}t)K\textbar\textbar{} ≤ δ\_eq +
βΔ\_t. Unroll the recursion. □

\hypertarget{block-contraction-for-sparse-kernels-and-the-w-t_mix-rule}{%
\subsection{9.2 Block contraction for sparse kernels and the w ≳ t\_mix
rule}\label{block-contraction-for-sparse-kernels-and-the-w-t_mix-rule}}

For sparse local kernels, one can have β(K)=1 because different rows can
have disjoint supports, making Lemma 9.1 inapplicable. One remedy is a
block coefficient τ\_m:= β(K\^{}m), which becomes \textless1 once m
exceeds the chain's mixing scale. \textbf{Lemma 9.2 (block-time
geometric bound).} Let δ\_m:= sup\_x TV(K\_c\^{}m(x,·), K\^{}m(x,·)).
Then for any k≥1, \textbar\textbar{} μ K\_c\^{}\{km\} − μ K\^{}\{km\}
\textbar\textbar\_TV ≤ δ\_m (1−τ\_m\^{}k)/(1−τ\_m) ≤ δ\_m/(1−τ\_m). A
crude but general bound is δ\_m ≤ m δ\_eq (telescoping), so the
stability term is O(m δ\_eq/(1−τ\_m)). To pick m, relate τ\_m to the
mixing profile d\_K(t):=max\_x TV(K\^{}t(x,·), π). \textbf{Lemma 9.3
(τ\_m from mixing).} For all t≥1, β(K\^{}t) ≤ 2 d\_K(t). In particular,
letting m=t\_mix(K,1/4) (so d\_K(m)≤1/4) gives τ\_m≤1/2. \textbf{Design
consequence (mixing-time preservation).} In the MC-EQ regime
δ\_eq=Θ(1/w), and taking m≈t\_mix(K) so that τ\_m is bounded away from
1, Lemma 9.2 suggests that keeping the realized process within O(ε) TV
of the cover process on the natural time scale t≈t\_mix(K) typically
requires w ≳ t\_mix(K)/ε. This formalizes the heuristic that w is
simultaneously a privacy knob (Sections 6-7) and a stability knob for
the cover process itself.

\hypertarget{flagship-compiler-calculus-theorem-composing-psc-mc-eq-and-stop-time-padding}{%
\section{10. Flagship compiler calculus theorem: composing PSC, MC-EQ,
and stop-time
padding}\label{flagship-compiler-calculus-theorem-composing-psc-mc-eq-and-stop-time-padding}}

We now state a single ``plug-and-play'' theorem that composes the main
compiler modules in this draft. It has two variants: (i) a
\textbf{posterior control} variant (near-\(1/K\) guarantees) based on
Theorem 7.2, and (ii) an \textbf{MI/Fano} variant (high-confidence
identification control) based on the Markov-reference theorem.

\hypertarget{posterior-control-variant-near-1k}{%
\subsection{\texorpdfstring{10.1 Posterior-control variant
(near-\(1/K\))}{10.1 Posterior-control variant (near-1/K)}}\label{posterior-control-variant-near-1k}}

\textbf{Theorem 10.1 (Compiler calculus: sequential posterior control
for joint traces).} Let \(C\sim\mathrm{Unif}([K])\) be a hidden profile.
Consider a joint transcript \((O_{1:T},\widetilde T)\) consisting of a
contact-mark trace \(O_{1:T}\) and a padded stop time \(\widetilde T\).
Assume: 1. (\textbf{PSC / public schedule}) The set of admissible
actions at time \(t\) is a public set \(J_t\) of size \(m_t\), and the
reference distribution \(K_t(\cdot\mid h)\) supported on \(J_t\)
satisfies a lower bound \(\kappa_t:= \inf_{h,y} K_t(y\mid h) > 0\). 2.
(\textbf{MC-EQ / per-history additive deviation}) For each \(t\le T\),
for all profiles \(c\), histories \(h\), and outcomes \(y\), \[
|P(O_t=y\mid C=c,H_{t-1}=h) - K_t(y\mid h)| \le \delta_t < \kappa_t.
\] A sufficient condition is the one-real-among-\(w\) mixing form, in
which Lemma 4.2 gives the pointwise deviation bound \(\delta_t\le 1/w\)
(and Corollary 4.1 gives the TV deviation \(\delta_{eq,t}\le 1/w\)). 3.
(\textbf{Stop-time padding with a posterior-friendly reference}) There
exists a secret-independent reference distribution \(R\) on
\(\{1,\dots,n\}\) with \(\kappa_{T+1}:=\min_t R(t)>0\) such that for all
\(c,h,y\), \[
|P(\widetilde T=y\mid C=c,H_T=h)-R(y)|\le \delta_{T+1} < \kappa_{T+1}.
\] For example, this holds if one solves PosteriorStopTime-Pad (Problem
2.14) with parameters \((\kappa_T,\delta_T)\). Then for every realized
transcript \((y_{1:T},t)\), \[
\max_{c} \Pr[C=c\mid O_{1:T}=y_{1:T},\widetilde T=t] \le \frac{\exp(\varepsilon_{T+1})}{K},
\qquad
\varepsilon_{T+1}:= \sum_{t=1}^{T+1} \log\!\left(\frac{\kappa_t+\delta_t}{\kappa_t-\delta_t}\right).
\] In particular, to guarantee near-random identification
\(\max_c \Pr[C=c\mid \text{transcript}]\le (1+\eta)/K\), it suffices
that \(\varepsilon_{T+1}\le \log(1+\eta)\). \emph{Proof.} This is
exactly Theorem 7.2 applied to the length-\((T{+}1)\) transcript.
Assumptions (1)-(3) are the per-history additive deviation hypotheses
with a positive lower bound. \textbf{Design plug-in (uniform scheduled
sets).} If \(K_t\) is uniform on \(J_t\) (so \(\kappa_t=1/m_t\)) and
one-real-among-\(w\) mixing gives \(\delta_t\le 1/w\) (more sharply,
\(\delta_t=(1-1/m_t)/w\) in the canonical injection), then \[
\varepsilon_{T+1}\le \sum_{t=1}^T \log\!\left(\frac{1/m_t+1/w}{1/m_t-1/w}\right) + \log\!\left(\frac{\kappa_T+\delta_T}{\kappa_T-\delta_T}\right).
\] When \(m_t/w\) is small this is approximately \(\sum_t 2m_t/w\) plus
the stop-time term (using \(\log\frac{1+a}{1-a}=2a+O(a^3)\) with
\(a=m_t/w\)).

\hypertarget{mifano-variant-high-confidence-identification}{%
\subsection{10.2 MI/Fano variant (high-confidence
identification)}\label{mifano-variant-high-confidence-identification}}

\textbf{Theorem 10.2 (Compiler calculus: MI bound + Fano inversion).}
Assume the Markov-reference hypothesis of Section 6 holds for
\(O_{1:T}\) with history-uniform TV deviation \(\delta_{eq}\) and
reference kernel minimum \(k_{\min}>0\). Assume also that the stop-time
compiler satisfies a per-history TV deviation bound
\(\mathrm{TV}(P(\widetilde T\mid C=c,H_T=h), R(\cdot\mid h))\le \delta_T\)
for some secret-independent reference \(R(\cdot\mid h)\) with
\(\kappa_T:=\inf_{h,t}R(t\mid h)>0\). Then \[
I(C;O_{1:T},\widetilde T) \le (T-1)\ln\!\left(1+\frac{2\delta_{eq}^2}{k_{\min}}\right) + \ln\!\left(1+\frac{2\delta_T^2}{\kappa_T}\right) + I(C;O_1).
\] Consequently, if \(C\) is uniform on \([K]\), Fano gives
\(P_{\mathrm{succ}}\le (I+\ln 2)/\ln K\). \textbf{Lemma 10.3 (KL from TV
and a minimum mass).} Let \(P,Q\) be distributions on a finite set with
\(q_{\min}:=\min_y Q(y)\ge \kappa>0\) and
\(\mathrm{TV}(P,Q)\le \delta<\kappa\). Then \[
D_{\mathrm{KL}}(P\|Q)\le \ln\!\left(1+\frac{2\,\delta^2}{\kappa}\right).
\] \emph{Proof.} Let \(\Delta(y):=P(y)-Q(y)\) and write
\(\Delta_+(y):=\max\{\Delta(y),0\}\) and
\(\Delta_-(y):=\max\{-\Delta(y),0\}\). Then
\(\sum_y \Delta_+(y)=\sum_y \Delta_-(y)=\mathrm{TV}(P,Q)=\delta\) and \[
\sum_y \Delta(y)^2 = \sum_y \Delta_+(y)^2 + \sum_y \Delta_-(y)^2 \le \delta^2+\delta^2 = 2\delta^2.
\] Hence the chi-square divergence satisfies \[
\chi^2(P\|Q):=\sum_y \frac{(P(y)-Q(y))^2}{Q(y)} \le \frac{1}{\kappa}\sum_y \Delta(y)^2 \le \frac{2\delta^2}{\kappa}.
\] Finally, by Jensen's inequality, \[
D_{\mathrm{KL}}(P\|Q)=\mathbb{E}_P[\ln(P/Q)]\le \ln\mathbb{E}_P[P/Q]
= \ln\sum_y \frac{P(y)^2}{Q(y)} = \ln\!\bigl(1+\chi^2(P\|Q)\bigr),
\] which yields the claim. \(\square\) \emph{Proof of Theorem 10.2.} By
the Markov-reference theorem (Section 6), \[
I(C;O_{1:T})\le (T-1)\ln\!\left(1+\frac{2\delta_{eq}^2}{k_{\min}}\right)+ I(C;O_1).
\] For the padded stop time, condition on the history \(H_T=h\). Let
\(P_{c,h}\) denote the conditional law of \(\widetilde T\) given
\((C=c,H_T=h)\), and let \(\bar P_h:=\frac1K\sum_c P_{c,h}\) be the
profile mixture. Then \[
I(C;\widetilde T\mid H_T=h)=\frac1K\sum_c D_{\mathrm{KL}}(P_{c,h}\|\bar P_h)\le \frac1K\sum_c D_{\mathrm{KL}}(P_{c,h}\|R_h),
\] since
\(\sum_c D(P_{c,h}\|\bar P_h)=\sum_c D(P_{c,h}\|R_h)-K\,D(\bar P_h\|R_h)\le \sum_c D(P_{c,h}\|R_h)\).
By the assumed per-history TV bound and Lemma 10.3, \[
D_{\mathrm{KL}}(P_{c,h}\|R_h)\le \ln\!\left(1+\frac{2\delta_T^2}{\kappa_T}\right)
\] for every \(c,h\). Thus
\(I(C;\widetilde T\mid H_T)\le \ln(1+2\delta_T^2/\kappa_T)\). Finally,
by the chain rule, \[
I(C;O_{1:T},\widetilde T)=I(C;O_{1:T})+I(C;\widetilde T\mid O_{1:T})
\le I(C;O_{1:T})+I(C;\widetilde T\mid H_T),
\] which gives the stated bound. \(\square\)

\hypertarget{a-small-sanity-check-non-adversarial}{%
\subsection{10.3 A small sanity check
(non-adversarial)}\label{a-small-sanity-check-non-adversarial}}

To gauge constant factors, we simulated the simplest i.i.d. model
consistent with Theorem 10.1: \(K=200\) profiles, a uniform scheduled
set of size \(m=20\), horizon \(T=20\), one-real-among-\(w\) injection
with \(w=2000\), and a stop-time mixture with \(n=10\) and
\(\lambda=0.05\). Across 2,000 trials, the empirical median of
\(\max_c \Pr[C=c\mid \text{transcript}]\) was about \(0.00725\), while
the Theorem 7.2 bound predicted \(\le 0.0193\). Meanwhile, an MI/Fano
accounting bound was numerically much looser in this near-\(1/K\)
regime. See \texttt{compiler\_calculus\_sims\_v23.py}.

\hypertarget{design-worksheet-worked-example}{%
\subsection{10.4 Design worksheet (worked
example)}\label{design-worksheet-worked-example}}

This subsection shows how to instantiate Theorem 10.1 with concrete
numbers. We emphasize that the resulting bounds are \emph{worst-case}
and can be conservative; the goal is a simple, checkable knob-setting
recipe. \textbf{Inputs.} Choose: - number of profiles \(K\) (uniform
prior), - trace horizon \(T\), - public scheduled set sizes
\(m_t:=|J_t|\), - desired posterior slack \(\eta\) (target
\(\max_c \Pr[C=c\mid \text{transcript}]\le (1+\eta)/K\)), - stop-time
reference support size \(n\) and a stop-time additive deviation target
\((\kappa_{T+1},\delta_{T+1})\). \textbf{Step 1 (mark-trace term).} If
each step uses one-real-among-\(w\) mixing inside a uniform scheduled
set of size \(m_t\) and \(w>m_t\), then \(\kappa_t=1/m_t\) and Lemma 4.2
gives a per-step pointwise deviation \(\delta_t\le 1/w\) (sharper:
\(\delta_t=(1-1/m_t)/w\) in the canonical injection). Compute \[
\varepsilon_{\mathrm{trace}}(w):=\sum_{t=1}^T \ln\!\left(\frac{1/m_t+1/w}{1/m_t-1/w}\right).
\] \textbf{Step 2 (stop-time term).} Choose (or optimize for) a
secret-independent stop-time reference \(R\) on \(\{1,\dots,n\}\) with
\(\kappa_{T+1}:=\min_t R(t)\). One can enforce the needed additive
deviation and floor constraints by solving PosteriorStopTime-Pad
(Problem 2.14), and then ensure
\(|P(\widetilde T=t\mid c,h)-R(t)|\le \delta_{T+1}<\kappa_{T+1}\) for
all \(c,h,t\). For example, if \(R\) is uniform then
\(\kappa_{T+1}=1/n\), and a conservative sufficient condition under a
simple mixture stop-time compiler is \(\delta_{T+1}\le \lambda\). A
convenient way to enforce these pointwise and floor constraints is to
solve \textbf{PosteriorStopTime-Pad} (Problem 2.14) for the desired
\((\kappa_{T+1},\delta_{T+1})\); any feasible padding kernel from that
LP plugs directly into Theorem 10.1 via Proposition 2.15. \textbf{Step 3
(budgeting).} Theorem 10.1 yields \[
\max_c \Pr[C=c\mid \text{transcript}] \le \frac{\exp(\varepsilon_{\mathrm{trace}}(w)+\varepsilon_{\mathrm{stop}})}{K},
\qquad
\varepsilon_{\mathrm{stop}}:=\ln\!\left(\frac{\kappa_{T+1}+\delta_{T+1}}{\kappa_{T+1}-\delta_{T+1}}\right).
\] Thus it suffices to pick \(w\) and stop-time parameters so that
\(\varepsilon_{\mathrm{trace}}(w)+\varepsilon_{\mathrm{stop}}\le \ln(1+\eta)\).
\textbf{Worked examples.} Take \(K=10^5\), \(T=50\), and uniform
scheduled sets of size \(m_t\equiv 20\). For stop time, take \(n=30\)
with uniform reference \(R\), so \(\kappa_{T+1}=1/30\), and target
\(\delta_{T+1}=0.001\), which gives
\(\varepsilon_{\mathrm{stop}}=\ln((1/30+0.001)/(1/30-0.001))\approx 0.0600\).
(A) \textbf{Near-random guessing within a factor 2.} Aim for
\(\max_c \Pr[C=c\mid \text{transcript}]\le 2/K\) (i.e.~\(\eta=1\), so
\(\ln(1+\eta)=\ln 2\)). This leaves an allowable trace budget of about
\(0.6331\). Solving \(\varepsilon_{\mathrm{trace}}(w)\le 0.6331\) yields
\(w=3159\) as the smallest integer meeting the bound: \[
\varepsilon_{\mathrm{trace}}(3159)=50\ln\left(\frac{1/20+1/3159}{1/20-1/3159}\right)\approx 0.63312,
\qquad\varepsilon_{\mathrm{trace}}+\varepsilon_{\mathrm{stop}}\approx 0.69314\le\ln 2.
\] (B) \textbf{Within a factor 10.} Aim for
\(\max_c \Pr[C=c\mid \text{transcript}]\le 10/K\) (i.e.~\(\eta=9\),
\(\ln(1+\eta)=\ln 10\)). Then the smallest feasible \(w\) is \(w=892\):
\[
\varepsilon_{\mathrm{trace}}(892)\approx 2.24253,
\qquad\varepsilon_{\mathrm{trace}}+\varepsilon_{\mathrm{stop}}\approx 2.30253\le\ln 10.
\] (C) \textbf{Within a factor 100.} Aim for
\(\max_c \Pr[C=c\mid \text{transcript}]\le 100/K\) (i.e.~\(\eta=99\),
\(\ln(1+\eta)=\ln 100\)). Then the smallest feasible \(w\) is \(w=441\).

\begin{enumerate}
\def\labelenumi{(\Alph{enumi})}
\setcounter{enumi}{3}
\tightlist
\item
  \textbf{Non-uniform public schedules.} Consider \(T=100\) slots split
  into two public schedule sizes: \(50\) slots with \(m=8\) and \(50\)
  slots with \(m=32\). Take \(n=60\) and \(\delta_{T+1}=0.001\), so \[
  \varepsilon_{\mathrm{stop}} = \ln\left(\frac{1/60+0.001}{1/60-0.001}\right) \approx 0.12014.
  \] For posterior targets (uniform prior over \(K\) profiles):
\end{enumerate}

\begin{itemize}
\tightlist
\item
  \textbf{Within a factor 2.} For \(\eta=1\) (target \(\le 2/K\)), the
  minimal feasible \(w\) under the loose per-step ratio is \(w=6981\).
\item
  \textbf{Within a factor 10.} For \(\eta=9\) (target \(\le 10/K\)), the
  minimal feasible \(w\) is \(w=1833\).
\item
  \textbf{Within a factor 100.} For \(\eta=99\) (target \(\le 100/K\)),
  the minimal feasible \(w\) is \(w=893\).
\end{itemize}

Under the canonical injection ratio (uniform cover on the scheduled
set), these improve to \(w=3478,\ 904,\ 434\) respectively. All values
are reproduced by
\texttt{compiler\_design\_worksheet\_latest.py\ -\/-examples}.

These examples illustrate that \textbf{posterior-first} guarantees
become expensive when the goal is extremely close to random guessing
(small \(\eta\)). \textbf{Contrast (MI/Fano route).} If instead one only
seeks a high-confidence identification bound,
e.g.~\(P_{\mathrm{succ}}\le 0.2\) among \(K=10^5\) profiles, then
Theorem 10.2 combined with the exact uniform per-step KL approximation
(Section 5.3) suggests a much smaller cover width (about \(w\approx 17\)
in the same \(T=50,m=20\) setting). This gap is expected: MI/Fano
controls \emph{average} information, while Theorem 10.1 controls
worst-case posteriors transcript-by-transcript. A short script that
computes the smallest feasible \(w\) for given
\((T,m_t,n,\delta_{T+1},\eta)\) is included as
\texttt{compiler\_design\_worksheet\_latest.py}.

\textbf{Figure 3 (worksheet scaling).} The posterior-first design
worksheet in Section 10.4 implies that the required equalization fanout
\(w\) scales roughly like \(\sum_t m_t\) divided by the allowable
posterior slack \(\ln(1+\eta)-\varepsilon_{\mathrm{stop}}\), up to the
exact nonlinearity \(\log\frac{1+a}{1-a}\). The plot below shows \(w\)
vs \(\eta\) for the two schedules used in the worked examples (uniform
\(m=20\), \(T=50\); and mixed \(m\in\{8,32\}\), \(T=100\)), comparing
the loose per-step ratio \(a_t=m_t/w\) to the canonical injection ratio
\(a_t=(m_t-1)/w\).

\begin{center}
\fbox{\begin{minipage}{0.93\linewidth}\small
\textbf{Text-only replacement for Fig.~3.} The plot is omitted (text-only submission). It visualized the design worksheet curve:
the minimal cover width $w$ needed to hit posterior target $(1+\eta)/K$ decreases with $\eta$, and is typically much larger under transcript-by-transcript posterior control than under MI/Fano. (To regenerate the original plot, run \texttt{\detokenize{make_figures_v67.py}}.)
\end{minipage}}
\end{center}

\hypertarget{related-work-non-exhaustive}{%
\section{11. Related work
(non-exhaustive)}\label{related-work-non-exhaustive}}

This paper is intentionally modular: we treat metadata leakage channels
(stop time and trace/marks) as \emph{interfaces} that export
quantitative contracts, so that compilers can be analyzed by composition
rather than by a monolithic, system-level argument. Closest in spirit
are works that (i) give operationally meaningful leakage measures, and
(ii) provide contraction/perturbation tools for sequential observation
processes.

Our analysis draws on (a) pointwise leakage measures that directly
control posteriors, (b) strong data-processing and contraction
coefficients for Markovian observation processes, and (c) sharp KL--TV
inequalities used to turn TV-style compiler contracts into KL/MI bounds.
The following anchors touch the specific tools used here:

\begin{itemize}
\tightlist
\item
  \textbf{Pointwise leakage / maximal leakage.} Operational formulations
  and closed forms for maximal leakage and related pointwise measures
  appear in {[}IWM18{]} (which motivates our Section 7
  ``posterior-first'' stance).
\item
  \textbf{SDPIs and contraction coefficients.} Strong data-processing
  inequalities and contraction viewpoints aligned with our ``reference
  process + contraction'' accounting are developed in {[}PW15{]}.
\item
  \textbf{Doeblin/Dobrushin-type coefficients.} Structural relationships
  among Doeblin-style measures and contraction coefficients appear in
  {[}MS23{]}, supporting the Section 9 mixing/contraction bookkeeping.
\item
  \textbf{Perturbation sensitivity of Markov chains.} Quantitative
  stability bounds for uniformly ergodic chains under kernel
  perturbations appear in {[}Mit05{]}, conceptually close to our
  ``equalization as controlled perturbation'' lens.
\item
  \textbf{Reverse Pinsker and KL--TV envelopes.} Reverse Pinsker-type
  inequalities and related bounds appear in {[}Sas15{]}, while
  Bretagnolle--Huber {[}BH79{]} and a modern KL--TV note {[}Can22{]}
  motivate our Lemma 10.3 / Section 8 envelope discussion.
\item
  \textbf{Background information theory.} Standard identities and
  inequalities (chain rules, Fano, etc.) are in {[}CT06{]}.
\end{itemize}

\hypertarget{open-problems-updated}{%
\section{12. Open problems (updated)}\label{open-problems-updated}}

\begin{enumerate}
\def\labelenumi{\arabic{enumi}.}
\item
  \textbf{Global optimal stop-time padding beyond mixtures.}
  Kernel-TV-Pad gives the true optimum among kernel-feasible paddings,
  and we exhibit separating families. What is the best \emph{structural}
  characterization of optimal kernels (e.g., sparsity patterns,
  monotonicity, bounds on the number of randomized rows)? Concrete
  progress would include: (i) a theorem that mixtures achieve a
  constant-factor approximation to the LP for all instances; or (ii) an
  efficiently checkable condition under which the LP optimum \emph{is} a
  mixture; or (iii) lower-bound constructions that force multi-row
  randomization in a quantified way.
\item
  \textbf{Deriving history-uniform contracts from real retry/repair
  policies.} The compiler calculus consumes per-history TV (or
  pointwise) deviation bounds. Can we systematically derive such bounds
  from concrete overlay protocols (timeouts, retries, backoff, batching,
  repair), especially when the cover kernel is state-dependent? Progress
  would mean reusable lemmas of the form: ``Policy X implies
  \(\mathrm{TV}(P_t(\cdot\mid c,h),K_t(\cdot\mid h))\le \delta_t\) for
  all histories \(h\)'' with explicit \(\delta_t\) as a function of
  protocol parameters.
\item
  \textbf{Moment-based sequential KL/MI accounting without global
  floors.} Lemma 5.2 avoids a global \(q_{\min}\) for one-step KL via
  moments. Can analogous moment-style bounds propagate through
  sequential composition (Theorem 6.1 / 10.2) to avoid global
  \(k_{\min}\) assumptions, especially when observation alphabets are
  large and supports drift with history? Progress would mean MI/KL
  bounds stated in terms of per-step moments (or collision
  probabilities) that remain non-vacuous at scale.
\item
  \textbf{Tightness of near-\(1/K\) posterior bounds under scheduling +
  injection.} The sequential PML/posterior theorems yield clean
  sufficient conditions, but can be loose. When are the bounds tight (up
  to constants), and when do they blow up unnecessarily? Progress would
  include matching lower bounds (adversary strategies) for
  scheduled-contact mechanisms, or identifying symmetric regimes where
  an exact posterior analysis is possible.
\end{enumerate}

\hypertarget{references}{%
\section{References}\label{references}}

{[}IWM18{]} Ibrahim Issa, Aaron B. Wagner, and Sudeep Kamath. \emph{An
Operational Approach to Information Leakage}. IEEE Transactions on
Information Theory, 66(3):1625--1657, March 2020. Preprint:
arXiv:1807.07878 (2018).

{[}Mit05{]} A. Yu. Mitrophanov. \emph{Sensitivity and convergence of
uniformly ergodic Markov chains}. Journal of Applied Probability,
42(4):1003--1014, December 2005. DOI: 10.1239/jap/1134587812.

{[}PW15{]} Yury Polyanskiy and Yihong Wu. \emph{Strong data-processing
inequalities for channels and Bayesian networks}. arXiv:1508.06025
(2015).

{[}MS23{]} Anuran Makur and Japneet Singh. \emph{Doeblin Coefficients
and Related Measures}. arXiv:2309.08475 (2023). Journal-ref: IEEE
Transactions on Information Theory, 70(7), July 2024.

{[}Sas15{]} Igal Sason. \emph{On Reverse Pinsker Inequalities}.
arXiv:1503.07118 (2015).

{[}Can22{]} Clément L. Canonne. \emph{A short note on an inequality
between KL and TV}. arXiv:2202.07198 (2022).

{[}BH79{]} Jean Bretagnolle and Catherine Huber. \emph{Estimation des
densités: risque minimax}. Zeitschrift für Wahrscheinlichkeitstheorie
und verwandte Gebiete, 47:119--137, January 1979. DOI:
10.1007/BF00535278.

{[}CT06{]} Thomas M. Cover and Joy A. Thomas. \emph{Elements of
Information Theory} (2nd ed.). Wiley, 2006.

\hypertarget{appendix-a.-pipeline-diagram-and-module-pseudocode}{%
\section{Appendix A. Pipeline diagram and module
pseudocode}\label{appendix-a.-pipeline-diagram-and-module-pseudocode}}

This appendix provides the promised one-page \emph{compiler calculus}
pipeline diagram (Figure 1) and lightweight pseudocode for the core
modules. The intent is operational clarity: each box exports an explicit
contract that the composition theorems (Theorems 10.1--10.2) consume.
\textbf{Figure 1 (Compiler calculus pipeline).}

\begin{verbatim}
User/profile C Public time/schedule
 | |
 | (payload layer hides | PSC: public schedule J_t
 | contents; adversary | (action set per slot)
 | sees metadata) v
 | Allowed action set J_t
 | |
 | | MC-EQ: choose w actions within J_t
 | | (1 real + w-1 covers)
 v v
Real target action Observed mark O_t in J_t
 | |
 | | Contract exported:
 | | (TV) TV(P_t(·|c,h), K_t(·|h)) ≤ 1/w
 | | (L∞) |P_t(y|c,h)-K_t(y|h)| ≤ 1/w
 | v
 | Mark transcript O_{1:T}
 |
 | Base success time S_c
 | (secret-dependent)
 v
StopTimePad: kernel-feasible padding Φ(t|s)
 |
 | Output: padded stop time \tilde T
 | Contract exported (two options):
 | (TV) TV(P_c, P_{c'}) ≤ τ (Kernel-TV-Pad)
 | (post) additive deviation + floor vs reference R
 v
Joint transcript (O_{1:T}, \tilde T)
 |
 v
Composition theorem:
 - Theorem 10.1: posterior control (near 1/K)
 - Theorem 10.2: MI/Fano control
\end{verbatim}

\hypertarget{a.1-psc-public-scheduling}{%
\subsection{A.1 PSC: public
scheduling}\label{a.1-psc-public-scheduling}}

PSC publishes the action-set schedule \(J_t\) and enforces fixed-size
traffic to every shard in \(J_t\) each slot. This collapses shard-level
marks: the adversary learns \(J_t\) from public time, and only learns
the \emph{within-set} action identifier from \(O_t\).

\hypertarget{a.2-mc-eq-mixture-equalization-within-a-scheduled-set}{%
\subsection{A.2 MC-EQ: mixture equalization within a scheduled
set}\label{a.2-mc-eq-mixture-equalization-within-a-scheduled-set}}

Below is a reference implementation sketch for one slot.

\begin{verbatim}
Input: history h, public set J_t, cover kernel K_t(·|h) on J_t, w ≥ 1
1. Sample cover multiset COV of size w-1 i.i.d. from K_t(·|h)
2. Sample real action A_real from the client rule (may depend on secret and h)
3. Send one request to each action in {A_real} ∪ COV (duplicates allowed)
4. Output mark O_t as the multiset or the (canonicalized) contacted identifier
\end{verbatim}

\textbf{Contract.} For the induced distribution of the observable mark
\(O_t\) (after canonicalization), \(P_t(\cdot\mid C=c,H_{t-1}=h)\) is a
\(\rho=1/w\) injection mixture relative to \(K_t(\cdot\mid h)\), hence:
- \(\mathrm{TV}(P_t(\cdot\mid c,h), K_t(\cdot\mid h)) \le 1/w\), and -
\(|P_t(y\mid c,h)-K_t(y\mid h)|\le 1/w\) for all \(y\in J_t\) (Lemma
4.2). These are precisely the per-slot hypotheses consumed by Theorem
10.1.

\hypertarget{a.3-stoptimepad-kernel-feasible-padding}{%
\subsection{A.3 StopTimePad: kernel-feasible
padding}\label{a.3-stoptimepad-kernel-feasible-padding}}

A kernel-feasible stop-time padding is a lower-triangular stochastic
matrix \(\Phi(t\mid s)\) supported on \(t\ge s\).

\begin{verbatim}
Input: realized success time s, padding kernel Φ
1. Sample padded stop time \tilde T ~ Φ(·|s)
2. Run dummy traffic until time \tilde T and then terminate
Output: \tilde T
\end{verbatim}

Two useful exported contracts are: - \textbf{Pairwise-TV contract}
(Kernel-TV-Pad): for all secrets \(c,c'\),
\(\mathrm{TV}(P_c,P_{c'})\le \tau\). - \textbf{Posterior-friendly
contract} (PosteriorStopTime-Pad): existence of a reference \(R\) with
floor \(\kappa_T\) and per-time additive deviation \(\delta_T\), so that
\(|P_c(t)-R(t)|\le \delta_T\) and \(R(t)\ge \kappa_T\) for all \(t\).
This is a direct plug-in to Theorem 10.1 (via Proposition 2.15).

\hypertarget{a.4-how-the-contracts-compose}{%
\subsection{A.4 How the contracts
compose}\label{a.4-how-the-contracts-compose}}

Theorem 10.1 treats the stop time \(\tilde T\) as one more observation
step with reference \(R\). Theorem 10.2 instead converts per-step TV
deviations to per-step KL (Lemma 10.3) and sums KL via the chain rule
before applying Fano.

\hypertarget{appendix-b.-numerical-sanity-checks-non-adversarial}{%
\section{Appendix B. Numerical sanity checks
(non-adversarial)}\label{appendix-b.-numerical-sanity-checks-non-adversarial}}

These checks are \textbf{not} part of the formal proof. They are quick
Monte Carlo tests intended to catch algebra/constant mistakes and to
gauge looseness.

\begin{enumerate}
\def\labelenumi{\arabic{enumi})}
\item
  \textbf{Moment KL bound for one-real-among-\(w\) injection (Lemma
  5.2).}\\
  We sample a dimension \(m\) uniformly from \(\{2,\dots,40\}\), a base
  distribution \(q\) uniformly from the simplex (Dirichlet\((1)\)), an
  index \(j\) uniformly from \(\{1,\dots,m\}\), and \(w\) uniformly from
  \(\{2,\dots,50\}\). With \(\rho=1/w\) and \(p=(1-\rho)q+\rho e_j\), we
  check
  \[D_{\mathrm{KL}}(p\|q)\ \le\ \rho^2\Bigl(\frac{1}{q_j}-1\Bigr).\]
  Over 5000 trials we observed \textbf{no violations}. The worst
  observed ratio \(D_{\mathrm{KL}}/\text{bound}\) was about 0.595.
\item
  \textbf{KL→TV envelope numbers (Pinsker vs Bretagnolle--Huber).}\\
  The script prints a small table for
  \(\varepsilon\in\{0.01,0.03,0.05,0.1,0.2,0.3\}\) comparing Pinsker's
  \(2\varepsilon^2\) to Bretagnolle--Huber's \(-\ln(1-\varepsilon^2)\)
  and their ratio.
\end{enumerate}

The script used is \texttt{sanity\_sims\_v57.py} (invoked by the
one-command runner in Appendix C.3).

\hypertarget{appendix-c.-reproducibility-checklist}{%
\section{Appendix C. Reproducibility
checklist}\label{appendix-c.-reproducibility-checklist}}

This draft ships with small scripts that reproduce the main
numerical/certificate claims. All scripts are deterministic given their
hard-coded instances and should run in seconds on a laptop.

\hypertarget{c.1-stop-time-padding}{%
\subsection{C.1 Stop-time padding}\label{c.1-stop-time-padding}}

\begin{itemize}
\tightlist
\item
  \textbf{Binary separation (kernel-feasible beats mixtures)}:
  \texttt{stop\_time\_dual\_cert\_v16.py} (primal kernel + dual
  multipliers + KKT residuals).
\item
  \textbf{Multi-row optimal padding example}:
  \texttt{stop\_time\_multirow\_cert\_v47.py} (two-row randomization
  instance; HiGHS KKT residuals).
\item
  \textbf{Parametric family (closed form, linear-in-horizon gap)}:
  \texttt{stop\_time\_dual\_cert\_v18.py} (checks dual inequalities for
  the family and confirms the closed-form objective matches the LP).
\item
  \textbf{\(K\)-ary star vs pairwise gap (Proposition 2.12)}:
  \texttt{stop\_time\_star\_gap\_v21.py} (prints optimal
  kernels/objectives), and \texttt{stop\_time\_kary\_cert\_v38.py}
  (extracts HiGHS dual multipliers and checks KKT).
\item
  \textbf{PosteriorStopTime-Pad certificate (Problem 2.14)}:
  \texttt{stop\_time\_posterior\_cert\_v46.py} (HiGHS primal--dual gap
  and KKT residuals for a small instance).
\item
  \textbf{TV geometry (Lemma 2.13)}:
  \texttt{tv\_radius\_geometry\_v24.py}.
\end{itemize}

\hypertarget{c.2-trace-posterior-design-worksheet}{%
\subsection{C.2 Trace / posterior / design
worksheet}\label{c.2-trace-posterior-design-worksheet}}

\begin{itemize}
\tightlist
\item
  \textbf{Posterior-composition sanity simulation} (for Theorem 7.2 /
  10.1 style bounds): \texttt{compiler\_calculus\_sims\_v23.py}.
\item
  \textbf{Exact-vs-bounding KL accounting for injection}:
  \texttt{w\_compare\_bounds.py}.
\item
  \textbf{Design worksheet (solve for \(w\) under posterior targets and
  MI/Fano contrast)}: \texttt{compiler\_design\_worksheet\_latest.py}.
\end{itemize}

\hypertarget{c.3-one-command-runner}{%
\subsection{C.3 One-command runner}\label{c.3-one-command-runner}}

\begin{itemize}
\tightlist
\item
  A convenience runner is provided as \texttt{reproduce\_all\_v67.sh}.
\end{itemize}

\hypertarget{c.4-figures}{%
\subsection{C.4 Figures}\label{c.4-figures}}

\begin{itemize}
\tightlist
\item
  \textbf{Figure 2 (gap vs horizon)} and \textbf{Figure 3 (worksheet
  scaling)}: \texttt{make\_figures\_v67.py} (writes
  \texttt{fig\_gap\_vs\_n} and \texttt{fig\_w\_vs\_eta}).
\end{itemize}

\hypertarget{appendix-d.-bibtex-entries-copypaste}{%
\section{Appendix D. BibTeX entries
(copy/paste)}\label{appendix-d.-bibtex-entries-copypaste}}

The following BibTeX entries cover all citation tags used in the draft.
They are intended as a convenience for submission; feel free to replace
with your preferred bibliography style.

\begin{verbatim}
@book{CT06,
  author    = {Thomas M. Cover and Joy A. Thomas},
  title     = {Elements of Information Theory},
  edition   = {2},
  publisher = {Wiley-Interscience},
  year      = {2006},
  isbn      = {978-0471241959}
}

@misc{IWM18,
  author       = {Ibrahim Issa and Aaron B. Wagner and Sudeep Kamath},
  title        = {An Operational Approach to Information Leakage},
  year         = {2018},
  eprint       = {1807.07878},
  archivePrefix= {arXiv},
  primaryClass = {cs.IT}
}

@misc{PW15,
  author       = {Yury Polyanskiy and Yihong Wu},
  title        = {Strong data-processing inequalities for channels and Bayesian networks},
  year         = {2015},
  eprint       = {1508.06025},
  archivePrefix= {arXiv},
  primaryClass = {cs.IT}
}

@misc{MS23,
  author       = {Anuran Makur and Japneet Singh},
  title        = {Doeblin Coefficients and Related Measures},
  year         = {2023},
  eprint       = {2309.08475},
  archivePrefix= {arXiv},
  primaryClass = {cs.IT}
}

@article{Mit05,
  author  = {A. Yu. Mitrophanov},
  title   = {Sensitivity and convergence of uniformly ergodic {M}arkov chains},
  journal = {Journal of Applied Probability},
  volume  = {42},
  number  = {4},
  pages   = {1003--1014},
  year    = {2005},
  doi     = {10.1239/jap/1134587812}
}

@misc{Can22,
  author       = {Cl{\'e}ment L. Canonne},
  title        = {A short note on an inequality between {KL} and {TV}},
  year         = {2022},
  eprint       = {2202.07198},
  archivePrefix= {arXiv},
  primaryClass = {cs.IT}
}

@misc{Sas15,
  author       = {Igal Sason},
  title        = {On Reverse Pinsker Inequalities},
  year         = {2015},
  eprint       = {1503.07118},
  archivePrefix= {arXiv},
  primaryClass = {cs.IT}
}

@incollection{BH79,
  author    = {Jean Bretagnolle and Catherine Huber},
  title     = {Estimation des densit{\'e}s: risque minimax},
  booktitle = {S{\'e}minaire de Probabilit{\'e}s {XII} (Univ. Strasbourg, Strasbourg, 1976/1977)},
  series    = {Lecture Notes in Mathematics},
  volume    = {649},
  pages     = {342--363},
  publisher = {Springer},
  address   = {Berlin},
  year      = {1978}
}
\end{verbatim}

\end{document}
